\SourcePage{139}
\SGAsection{5.5}{分裂约化群的 Borel 子群}
\begin{SGAresult}{命题}{5.5.1}
设 $(G,T,M,R)$ 为一个 $S$ 上的\Term{split-group}{分裂群}。对 $R$ 的每个\Term{positive-root-system}{正根系} $R_+$，$H_{R_+}$（由 5.4.7 存在）是 $G$ 的一个\Term{borel-subgroup}{Borel 子群}；对 $R_+$ 上的任意次序，由 $G$ 中乘法诱导的\Term{morphism}{态射}
\[T\times_S\prod_{\alpha\in R_+}U_\alpha\longrightarrow G\]
是像为 $H_{R_+}$ 的\Term{closed-immersion}{闭浸入}。记 $B_{R_+}=H_{R_+}$。
\end{SGAresult}
\begin{proof}
根据\Term{borel-subgroup}{Borel 子群}的定义，可将 $S$ 换成\Term{algebraically-closed-field}{代数闭域}的谱来验证第一个断言。这时令 $B$ 为 $G$ 中包含 $T$、对应于\Term{positive-root-system}{正根系} $R_+$ 的\Term{borel-subgroup}{Borel 子群}（Bible，\S 10.4，命题 9）；$B$ 的\Term{lie-algebra}{Lie 代数}为 $\mathfrak g_{R_+}$，故由 5.3.5 有 $B=H_{R_+}$。
证明第二个断言：题述\Term{morphism}{态射}诱导一个\Term{open-immersion}{开浸入} $i:T\times_S\prod_{\alpha\in R_+}U_\alpha\to H_{R_+}$（5.4.4），而 $i$ 是\Term{surjectivity}{满射}（Bible，\S 15.1，命题 1 的推论 1）。
\end{proof}
\SourcePage{140}
\begin{SGAresult}{推论}{5.5.2}
在 $R_+$ 上任取一个次序，并对每个 $\alpha\in R_+$ 取 $X_\alpha\in\Gamma(S,\mathfrak g^\alpha)^\times$。设 $\alpha,\beta\in R_+$。对每对满足 $i\alpha+j\beta\in R$ 的 $(i,j)\in\mathbb N^*\times\mathbb N^*$，存在唯一的\Term{section}{截面}
\[C_{i,j,\alpha,\beta}\in\Gamma(S,\mathcal O_S),\]
使得对所有 $S'\to S$ 及 $x,y\in\mathbb G_a(S')$ 有
\[\begin{split}
\exp_\alpha(xX_\alpha)\exp_\beta(yX_\beta)\exp_\alpha(xX_\alpha)^{-1}
={}&\exp_\beta(yX_\beta)\\[-2pt]
&\cdot\prod_{\substack{i,j\in\mathbb N^*\\i\alpha+j\beta\in R}}
\exp_{i\alpha+j\beta}(C_{i,j,\alpha,\beta}x^iy^jX_{i\alpha+j\beta}).
\end{split}\]
\end{SGAresult}
\begin{proof}
若 $\alpha=\beta$，断言显然成立。故设 $\alpha\ne\beta$；由上述命题，存在唯一的\Term{morphism}{态射}
\[F_0:\mathbb G_{a,S}^2\to T,\qquad F_\gamma:\mathbb G_{a,S}^2\to\mathbb G_{a,S}\quad(\gamma\in R_+),\]
使得
\[\exp(xX_\alpha)\exp(yX_\beta)\exp(xX_\alpha)^{-1}
=F_0(x,y)\prod_{\gamma\in R_+}\exp(F_\gamma(x,y)X_\gamma).\]
设 $t\in T(S')$，$S'\to S$。对上式施以 $\operatorname{int}(t)$，立即得到
\[F_0(\alpha(t)x,\beta(t)y)=F_0(x,y),\tag{1}\]
\[F_\gamma(\alpha(t)x,\beta(t)y)=\gamma(t)F_\gamma(x,y).\tag{2}\]
因 $\alpha,\beta$ 是 $T$ 的两个在 $\mathbb Q$ 上\Term{linearly-independent}{线性无关}的\Term{character}{特征标}，照常由第一个关系推出 $F_0$ 为常值，故 $F_0(x,y)=e$。再写成
\[F_\gamma(x,y)=\sum a_{ij}x^iy^j,\qquad a_{ij}\in\Gamma(S,\mathcal O_S).\]
代入关系 (2)，比较两端多项式，得到
\[a_{ij}\bigl(\alpha(t)^i\beta(t)^j-\gamma(t)\bigr)=0.\]
若 $\gamma\ne i\alpha+j\beta$，则存在\Term{faithfully-flat}{忠实平坦}且\Term{quasi-compact}{拟紧}的 $S'\to S$ 及 $t\in T(S')$，使 $\alpha(t)^i\beta(t)^j-\gamma(t)=1$（第 XIX 讲 4.13）。故在 $S'$ 上有 $a_{ij}=0$，从而在 $S$ 上亦然。若 $\gamma=i\alpha+j\beta$，置 $a_{ij}=C_{i,j,\alpha,\beta}$。令 $x=0$（分别令 $y=0$），得到 $C_{0,1,\alpha,\beta}=1$（分别为 $C_{1,0,\alpha,\beta}=0$）。
\end{proof}
\textbf{注记 5.5.3。} 在 $y=0$ 处求导并与 5.4.9 比较，得到
\[C_{i,1,\alpha,\beta}=M_{\alpha,\beta,i}.\]
\begin{SGAresult}{推论}{5.5.4}
设 $S$ 为\Term{scheme}{概形}，$G$ 为 $S$ 上的\Term{reductive-group}{约化群}，$T$ 为 $G$ 的\Term{maximal-torus}{极大环面}，$\alpha\ne\beta$ 为 $G$ 关于 $T$ 的两个\Term{root}{根}，并且 $\alpha+\beta$ 在每条\Term{fiber}{纤维}上均非平凡。将所有 $i\alpha+j\beta$（$i,j\in\mathbb N^*$）任意排序。对每个满足 $i\alpha+j\beta\in R$ 的 $i,j\in\mathbb N^*$，存在唯一的 $\mathcal O_S$ 模\Term{morphism}{态射}
\[f_{\alpha,\beta,i,j}:(\mathfrak g^\alpha)^{\otimes i}\otimes(\mathfrak g^\beta)^{\otimes j}\longrightarrow\mathfrak g^{i\alpha+j\beta},\]
\SourcePage{141}
使对每个 $S'\to S$ 及 $X\in\mathbf W(\mathfrak g^\alpha)(S')$、$Y\in\mathbf W(\mathfrak g^s)(S')$ 均有（右端的 $\exp$ 按 1.2 的意义理解\footnote{编者注：即若 $\mathfrak g^{i\alpha+j\beta}$ 在 $S$ 的某个\Term{connected-component}{连通分支}上为零，相应的指数在该分支上等于 $1$。}）
\[\exp_\alpha(X)\exp_\beta(Y)\exp_\alpha(-X)
=\exp_\beta(Y)\prod_{(i,j)}\exp_{i\alpha+j\beta}\bigl(f_{\alpha,\beta,i,j}(X^i\otimes Y^j)\bigr).\]
\end{SGAresult}
\begin{proof}
该断言对 (fpqc) 是\Term{local-scheme}{局部}的。由第 2 节，可设 $G$ 关于 $T$ \Term{split-group}{分裂}，且 $\alpha,\beta$ 在此\Term{splitting}{分裂}中为常值。因 $\alpha+\beta\ne0$，存在包含 $\alpha,\beta$ 的\Term{positive-root-system}{正根系} $R_+$（第 XXI 讲 3.5.4），于是归结为 5.5.2。
\end{proof}
\begin{SGAresult}{推论}{5.5.5}
设 $S$ 为\Term{scheme}{概形}，$G$ 为 $S$ 上的\Term{reductive-group}{约化群}。

(i) 对\Term{etale}{平展}\Term{topology}{拓扑}，$G$ \Term{local-scheme}{局部}具有\Term{borel-subgroup}{Borel 子群}。若 $T$ 为 $G$ 的\Term{maximal-torus}{极大环面}，则对\Term{etale}{平展}\Term{topology}{拓扑}，$G$ 也\Term{local-scheme}{局部}具有包含 $T$ 的\Term{borel-subgroup}{Borel 子群}。

(ii) 若 $T$ 为 $G$ 的\Term{maximal-torus}{极大环面}，则“$G$ 中包含 $T$ 的\Term{borel-subgroup}{Borel 子群}的\Term{functor}{函子}”可由 $W_G(T)$ 下的一个\Term{principal-homogeneous-space}{主齐性纤维丛}表示。

(iii) 若 $(G,T,M,R)$ \Term{split-group}{分裂}，则 $G$ 中包含 $T$ 的每个\Term{borel-subgroup}{Borel 子群} $B$，在 $S$ 上\Term{local-scheme}{局部}均形如 $B_{R_+}$，其中 $R_+$ 是 $R$ 的一个\Term{positive-root-system}{正根系}。

(iv) 若 $T\subset B$ 是 $G$ 的一个\Term{killing-pair}{Killing 对}，则存在\Term{etale}{平展}\Term{topology}{拓扑}的\Term{covering-family}{覆盖族} $\{S_i\to S\}$，并对每个 $i$ 存在一个\Term{splitting}{分裂} $(G_{S_i},T_{S_i},M_i,R_i)$ 及 $R_i$ 的一个\Term{positive-root-system}{正根系} $R_i^+$，使 $B_{S_i}=B_{R_i^+}$。
\end{SGAresult}
事实上，(i) 由 2.3 和 5.5.1 得到，(ii) 由 (i) 和 5.3.15 得到，(iii) 由 (ii) 和 5.5.1 得到，(iv) 由 (iii) 和 2.3 得到。
\begin{SGAresult}{引理}{5.5.6}
在由\Term{root}{根}生成的群 $\Gamma_0(R)$ 上取一个全序\Term{group-structure}{群结构}，使大于零的\Term{root}{根}恰为 $R_+$ 的元素（参见第 XXI 讲 3.5.6）\footnote{编者注：注意，这样的次序必与 $\operatorname{ord}_\Delta$ 相容，其中 $\Delta=S(R_+)$（参见 XXI 3.2.15）。}。令 $\alpha_1<\cdots<\alpha_N$ 为 $R_+$ 的所有元素。考虑由 $G$ 中乘法诱导的\Term{isomorphism}{同构}
\[f:T\times_SU_{\alpha_1}\times_S\cdots\times_SU_{\alpha_N}\longrightarrow B_{R_+}.\]
对 $i=1,\ldots,N$，置
\[U_{\ge i}=f(U_{\alpha_i}\times_S\cdots\times_SU_{\alpha_N}).\]
(i) 每个 $U_{\ge i}$ 均为 $B_{R_+}$ 的\Term{normal-subgroup}{正规子群}。

(ii) 对 $1\le i\le N-1$，$U_{\ge i}$ 与\Term{semidirect-product}{半直积}
\[U_{\ge i}=U_{\alpha_i}\cdot U_{\ge i+1}\]
等同。

(iii) $B_{R_+}$ 与\Term{semidirect-product}{半直积} $B_{R_+}=T\cdot U_{\ge1}$ 等同。
\SourcePage{142}
(iv) 对 $1\le i\le N-1$，$U_{\ge1}$ 的\Term{inner-automorphism}{内自同构}在 $U_{\ge i}/U_{\ge i+1}$ 上\Term{action}{作用}平凡（由 (ii)，此商与 $U_{\alpha_i}$ 等同）。
\end{SGAresult}
\begin{proof}
先对 $i$ 归纳证明：$U_{\ge i}$ 是 $B_{R_+}$ 的\Term{normal-subgroup}{正规子群}，且为 $U_{\alpha_i}$ 与 $U_{\ge i+1}$ 的\Term{semidirect-product}{半直积}。$i=N$ 时成立；设 $i+1$ 时成立，来证 $i$ 的情形。作为\Term{scheme}{概形}，有 $U_{\ge i}=U_{\alpha_i}\cdot U_{\ge i+1}$。先说明 $U_{\ge i}$ 在 $B_{R_+}$ 的\Term{inner-automorphism}{内自同构}下\Term{stable}{稳定}。这对 $\operatorname{int}(t)$（$t\in T(S)$）显然，只须对 $\operatorname{int}(x)$（$x\in U_\alpha(S)$、$\alpha\in R_+$）验证。因已假定 $U_{\ge i+1}$ \Term{normal}{正规}，只须说明 $\operatorname{int}(x)U_{\alpha_i}\subset U_{\ge i}$。由 5.5.2，若 $y\in U_{\alpha_i}(S')$，则 $y^{-1}xyx^{-1}\in U_{\ge i+1}(S')$，从而 $\operatorname{int}(x)(y)\in U_{\ge i}(S')$。
再证明 $U_{\ge i}$ 是 $B_{R_+}$ 的\Term{subgroup}{子群}。若 $x,y\in U_{\ge i}(S)$，可写成 $x=px'$、$y=qy'$，其中 $p,q\in U_{\alpha_i}(S)$，$x',y'\in U_{\ge i+1}(S)$。有
\[xy=px'qy'=pq(q^{-1}x'q)y'\in U_{\alpha_i}(S')U_{\ge i+1}(S');\]
同样，$x^{-1}=p^{-1}(px'^{-1}p^{-1})\in U_{\alpha_i}(S')U_{\ge i+1}(S')$。于是证明了 (i)、(ii)，并在过程中证明了 (iv)。(iii) 则直接由 5.5.1 得到。
\end{proof}
\begin{SGAresult}{引理}{5.5.7}
沿用上述记号，对每个 $1\le i\le N$ 取 $X_i\in\Gamma(S,\mathfrak g^{\alpha_i})^\times$，并考虑逐点由
\[a(x_1,\ldots,x_N)=\exp_{\alpha_1}(x_1X_1)\cdots\exp_{\alpha_N}(x_NX_N)\]
定义的\Term{isomorphism}{同构} $a:\mathbb G_{a,S}^N\to U_{\ge1}$。存在唯一一族系数在 $\Gamma(S,\mathcal O_S)$ 中的多项式
\[Q_i=Q_i(x_1,\ldots,x_N,y_1,\ldots,y_N)\qquad(i=1,\ldots,N),\]
使得逐点有
\[a(x_1,\ldots,x_N)a(y_1,\ldots,y_N)=a\bigl(Q_1(x_1,\ldots,y_N),\ldots,Q_N(x_1,\ldots,y_N)\bigr).\]
此外，$Q_i$ 的系数属于由 5.5.2 的 $C_{i,j,\alpha,\beta}$（$\alpha,\beta\in R_+$，$i,j\in\mathbb N^*$）生成的 $\Gamma(S,\mathcal O_S)$ 的子环，且每个 $Q_i$ 形如
\[Q_i(x_1,\ldots,y_N)=x_i+y_i+Q'_i(x_1,\ldots,x_{i-1},y_1,\ldots,y_{i-1}).\]
\end{SGAresult}
\begin{proof}
$Q_i$ 的存在唯一性立即由 $a$ 是\Term{scheme}{概形}\Term{isomorphism}{同构}得到。用 $z,z',z''$ 表示 $U_{\ge i+1}$ 的\Term{section}{截面}，有
\[\begin{split}
a(x_1,\ldots,x_N)a(y_1,\ldots,y_N)
={}&a(x_1,\ldots,x_{i-1},0,\ldots,0)\exp(x_iX_i)z\\
&\cdot a(y_1,\ldots,y_{i-1},0,\ldots,0)\exp(y_iX_i)z'.
\end{split}\]
利用 5.5.6 (i)、(iv)，右端可写成
\[a(x_1,\ldots,x_{i-1},0,\ldots,0)a(y_1,\ldots,y_{i-1},0,\ldots,0)\exp((x_i+y_i)X_i)z'';\]
再次利用 5.5.6，得到
\[Q_i(x_1,\ldots,x_N,y_1,\ldots,y_N)=x_i+y_i+Q'_i(x_1,\ldots,x_{i-1},y_1,\ldots,y_{i-1}),\]
\SourcePage{143}
其中
\[Q'_i(x_1,\ldots,x_{i-1},y_1,\ldots,y_{i-1})
=Q_i(x_i,\ldots,x_{i-1},0,\ldots,0,y_1,\ldots,y_{i-1},\ldots,0),\]
这就是所求的具体\Term{form}{形式}。
最后证明关于 $Q_i$ 系数的断言。令 $A$ 为由 $C_{i,j,\alpha,\beta}$（$\alpha,\beta\in R_+$，$i,j\in\mathbb N^*$）生成的 $\Gamma(S,\mathcal O_S)$ 的子环。对 $i$ 作递降归纳，证明若 $x_1=\cdots=x_{i-1}=0$ 且 $y_1=\cdots=y_{i-1}=0$，即 $a(x_1,\ldots,x_N)$、$a(y_1,\ldots,y_N)$ 是 $U_{\ge i}$ 的\Term{section}{截面}，则多项式
\[R_j(x_i,\ldots,x_N,y_i,\ldots,y_N)=Q_j(x_1,\ldots,x_N,y_1,\ldots,y_N)\]
的系数在 $A$ 中。这对 $i=N$ 显然，对 $j<i$ 也显然（此时 $R_j=0$）。设 $i<N$，断言对 $i+1$ 成立，来证 $i$（及 $j\ge i$）的情形。有
\[a(0,\ldots,0,x_i,\ldots,x_N)=\exp(x_iX_i)a(0,\ldots,0,x_{i+1},\ldots,x_N)=\exp(x_iX_i)Z_i.\]
同样写成 $a(0,\ldots,y_i,\ldots,y_N)=\exp(y_iX_i)T_i$，便有
\[\begin{split}
a(0,\ldots,x_i,\ldots,x_N)a(0,\ldots,y_i,\ldots,y_N)
=\exp((x_i+y_i)X_i)\operatorname{int}(\exp(-y_iX_i))(Z_i)\cdot T_i.
\end{split}\]
而
\[\operatorname{int}(\exp(-y_iX_i))(Z_i)
=\operatorname{int}(\exp(-y_iY_i))\bigl(\exp(x_{i+1}X_{i+1})\cdots\exp(x_NX_N)\bigr)\]
是 $U_{\ge i+1}$ 的 $N-i-1$ 个\Term{section}{截面}的乘积；由 5.5.2，这些\Term{section}{截面}在\Term{resolution}{分解} $U_{\ge i+1}=U_{\alpha_{i+1}}\cdots U_{\alpha_N}$ 中的坐标是 $y_i,x_{i+1},\ldots,x_N$ 的系数在 $A$ 中的多项式。应用归纳假设，推出
\[\operatorname{int}(\exp(-y_iX_i))(Z_i)\cdot T_i\]
的坐标也是系数在 $A$ 中的多项式，证明完成。
\end{proof}
注意，上述归纳立即给出下述引理的证明。
\begin{SGAresult}{引理}{5.5.8}
沿用 5.5.6 的记号，对每个 $i=1,\ldots,N$，设有群\Term{morphism}{态射} $f_i:U_{\alpha_i}\to H$，其中 $H$ 为一个 $S$ 上的群值\Term{functor}{函子}。由
\[f\bigl(\exp(x_1X_1)\cdots\exp(x_NX_N)\bigr)
=f_1(\exp(x_1X_1))\cdots f_N(\exp(x_NX_N))\]
定义的\Term{morphism}{态射} $f:U_{\ge1}\to H$ 是群\Term{morphism}{态射}，当且仅当对每对 $i<j$ 有
\[\begin{split}
f_j(\exp(x_jX_j))f_i(\exp(x_iX_i))f_j(\exp(-x_jX_j))
=f\bigl(\exp(x_jX_j)\exp(x_iX_i)\exp(-x_jX_j)\bigr).
\end{split}\]
\end{SGAresult}
\SourcePage{144}
\SGAsection{5.6}{具有可解纤维的 (R) 型子群}
\begin{SGAresult}{命题}{5.6.1}
设 $(G,T,M,R)$ 为一个 $S$ 上的\Term{split-group}{分裂群}，$R'$ 为 $R$ 的一个 (R) 型\Term{subset}{子集}（5.4.2），$H_{R'}$ 为 $G$ 的相应\Term{subgroup}{子群}。下列条件\Term{equivalence}{等价}：

(i) $H_{R'}$ 的\Term{geometric-fibre}{几何纤维}\Term{solvable}{可解}。

(ii) 存在一个\Term{positive-root-system}{正根系} $R_+$ 使 $R'\subset R_+$，从而 $H_{R'}\subset B_{R_+}$（参见 5.4.5）。

(iii) $R'\cap -R'=\varnothing$。

(iv) 对 $R'$ 上的任意次序，由 $G$ 中乘法诱导的\Term{morphism}{态射}
\[T\times_S\prod_{\alpha\in R'}U_\alpha\longrightarrow H_{R'}\]
是\Term{isomorphism}{同构}。

(v) $H_{R'}\cap\operatorname{Norm}_G(T)=T$。

(vi) 对 $R$ 的任意 (R) 型\Term{subset}{子集} $R''$，有（参见 5.4.5）
\[H_{R'}\cap\operatorname{Norm}_G(H_{R''})=H_{R'\cap R''}.\]
\end{SGAresult}
\begin{proof}
我们按照下面的逻辑图证明这些\Term{equivalence}{等价}性：
\[\begin{tikzcd}[column sep=large,row sep=small]
\text{(iii)} & \text{(ii)}\arrow[l,Rightarrow]\arrow[r,Rightarrow]\arrow[d,Rightarrow] & \text{(vi)}\arrow[dd,Rightarrow]\\
&\text{(iv)}\arrow[dr,Rightarrow]&\\
\text{(i)}\arrow[uu,Leftarrow]\arrow[uur,Rightarrow]&&\text{(v)}\arrow[ll,Rightarrow]
\end{tikzcd}\]
显然 (ii)$\Rightarrow$(iii)，且 (vi)$\Rightarrow$(v)（取 $R''=\varnothing$）。由 5.4.6，只须在\Term{geometric-fibre}{几何纤维}上验证 (i)$\Rightarrow$(ii)；若 $S$ 是\Term{algebraically-closed-field}{代数闭域}的谱，$H_{R'}$ 包含于一个含 $T$ 的\Term{borel-subgroup}{Borel 子群}中，此\Term{subgroup}{子群}因而形如 $H_{R_+}$（5.5.5 (iii)）。
同样，(iii)$\Rightarrow$(i) 可在\Term{geometric-fibre}{几何纤维}上验证。设 (iii) 成立；若 $H_{R'}$ 不\Term{solvable}{可解}，则 $T$ 中存在一个\Term{codimension}{余维}为 $1$ 的\Term{subtorus}{子环面} $Q$，使 $\operatorname{Centr}_{H_{R'}}(Q)$ 不\Term{solvable}{可解}（Bible，\S 10.4，命题 8）。而 $\operatorname{Centr}_{H_{R'}}(Q)$ 的\Term{lie-algebra}{Lie 代数}是 $\mathfrak g_{R''}$，其中 $R''$ 为 $R'$ 中在 $Q$ 上为零的\Term{root}{根}的\Term{set}{集合}；由 (iii)，$R''=\varnothing$ 或 $\{\alpha\}$。因此 $\operatorname{Centr}_{H_{R'}}(Q)$ 作为 $G$ 的一个 (R) 型\Term{subgroup}{子群}，等于 $T$ 或 $T\cdot U_\alpha$，故\Term{solvable}{可解}，与假设矛盾。
同样，(ii)$\Rightarrow$(iv) 可在\Term{geometric-fibre}{几何纤维}上验证（因所涉及的 $S$ 上\Term{scheme}{概形}均\Term{flat}{平坦}且\Term{finite-presentation}{有限呈示}）；由 Bible，\S 13.2，定理 1 d)，题述\Term{morphism}{态射}是\Term{bijection}{双射}；它在原点的\Term{tangent-space}{切空间}上诱导\Term{isomorphism}{同构}，照常便得结论（参见 4.1.1）。
由 4.2.7 有 (iv)$\Rightarrow$(v)。为证 (v)$\Rightarrow$(i)，再次归结到 $S$ 是\Term{algebraically-closed-field}{代数闭域}的谱的情形，再由 Bible，\S 10.3，命题 6 的推论以及 \S 9.3，定理 1 的推论 3 得出结论。
\SourcePage{145}
所以只须证明 (ii)$\Rightarrow$(vi)。可归结到 $G$ 是\Term{adjoint-group}{伴随群}的情形。由 5.3.3 有
\[\operatorname{Norm}_G(H_{R''})=\operatorname{Norm}_G(\mathfrak g_{R''})\subset\operatorname{Transp}_G(\mathfrak t,\mathfrak g_{R''}).\]\footnote{编者注：此包含关系由 $\mathfrak t\subset\mathfrak g_{R''}$ 得到。}
由 5.4.5，只须证明
\[H_{R'}(S)\cap\operatorname{Transp}_{G(S)}(\mathfrak t,\mathfrak g_{R''})\subset H_{R'\cap R''}(S).\tag{x}\]
先证明一个引理。
\end{proof}
\begin{SGAresult}{引理}{5.6.2}
沿用 5.5.7 的记号，设 $u=\exp(x_1X_1)\cdots\exp(x_NX_N)$，其中 $x_i\in\mathbb G_a(S)$。设 $m$ 为整数，$1\le m\le N$，且 $i<m$ 时 $x_i=0$。

a) 若 $H\in\Gamma(S,\mathfrak t)$，则 $\operatorname{Ad}(u)H$ 在 $\mathfrak g^{\alpha_m}$ 上的\Term{set-component}{分量}为 $-\overline\alpha_m(H)x_mX_m$。

b) 若 $Y\in\Gamma(S,\mathfrak g^{-\alpha_m})$，则按第 XX 讲 2.6 的记号，$\operatorname{Ad}(u)Y$ 在 $\mathfrak t$ 上的\Term{set-component}{分量}为 $x_m\langle X_m,Y\rangle H_{\alpha_m}$。
\end{SGAresult}
\begin{proof}
记 $\mathfrak u=\mathfrak g^{\alpha_{m+1}}+\cdots+\mathfrak g^{\alpha_N}$。由 5.4.9，$i>m$ 时有
\[\operatorname{Ad}(\exp(x_iX_i))\mathfrak u\subset\mathfrak u.\]
由第 XX 讲 2.10，有 $\operatorname{Ad}(\exp(x_iX_i))H=H-\overline\alpha_i(H)x_iX_i$。于是模 $\mathfrak u$ 立即有
\[\operatorname{Ad}(u)H=\operatorname{Ad}(\exp(x_mX_m))H=H-\overline\alpha_m(H)x_mX_m,\]
从而得到第一个结论。
同样，记\footnote{编者注：下文已对原版作了修正。} $\mathfrak n=\mathfrak g^{\alpha_1}+\cdots+\mathfrak g^{\alpha_N}$ 及 $u_1=\exp(x_{m+1}X_{m+1})\cdots\exp(x_NX_N)$。当 $i>m$ 时，$\alpha_i>\alpha_m$，故由 5.4.9，模 $\mathfrak n$ 有
\[\operatorname{Ad}(u_1)Y\equiv Y,\qquad\text{因而}\quad\operatorname{Ad}(u)Y\equiv\operatorname{Ad}(\exp(x_mX_m))Y.\]
应用第 XX 讲 2.10，模 $\mathfrak n$ 得
\[\operatorname{Ad}(u)Y-Y\equiv x_m\langle X_m,Y\rangle H_{\alpha_m},\]
从而得到第二个结论。
\end{proof}
\begin{proof}[继续证明 5.6.1]
回到包含关系 (x) 的证明。假设存在 $h\in H_{R'}(S)$、$h\notin H_{R'\cap R''}(S)$，使 $\operatorname{Ad}(h)\mathfrak t\subset\mathfrak g_{R''}$。可写成
\[h=t\exp(x_1X_1)\cdots\exp(x_NX_N).\]
\SourcePage{146}
由于 $h\notin H_{R'\cap R''}(S)$，存在最小的 $m$，使
\[t\exp(x_1X_1)\cdots\exp(x_{m-1}X_{m-1})\in H_{R'\cap R''}(S),\qquad\alpha_m\notin R'',\quad x_m\ne0.\]
于是 $h'=\exp(x_mX_m)\cdots\exp(x_NX_N)$ 也满足上述对 $h$ 的条件。但由 5.6.2，对每个 $H\in\Gamma(S,\mathfrak t)$，$\operatorname{Ad}(h')H$ 在 $\mathfrak g^{\alpha_m}$ 上的\Term{set-component}{分量}为 $-\overline\alpha_m(X)x_mX_m$。由对 $h$ 和 $m$ 的假设，对每个 $H\in\Gamma(S,\mathfrak t)$ 均有 $\overline\alpha_m(H)=0$，这是不可能的（因为 $G$ 假定为\Term{adjoint-group}{伴随群}，故 $\overline\alpha_m$ 在每条\Term{fiber}{纤维}上非零）。
\end{proof}
\textbf{注记 5.6.2 bis。} 沿用 5.6.2 的记号。若 $\operatorname{Ad}(u)$ 在 $\mathfrak t$ 和 $\mathfrak g^{-\alpha_m}$ 上为恒等，则 $x_m=0$。事实上，$x_m\overline\alpha_m=0$ 且 $x_mH_{\alpha_m}=0$；若 $\alpha_m\notin2M$，则 $\overline\alpha_m$ 在每条\Term{fiber}{纤维}上非零；若 $\alpha_m\in2M$，则 $\alpha_m^*\notin2M^*$，故 $H_{\alpha_m}$ 在每条\Term{fiber}{纤维}上非零；两种情形均推出 $x_m=0$。因此，若 $\operatorname{Ad}(u)$ 在 $\mathfrak g$ 上\Term{action}{作用}平凡，则 $u=e$。

\textbf{注记 5.6.3。} 若 $H$ 为\Term{reductive-group}{约化群} $G$ 的一个\Term{geometric-fibre}{几何纤维}\Term{solvable}{可解}的 (R) 型\Term{subgroup}{子群}，则 $H$ 在 $G$ 中闭，且 $\operatorname{Norm}_G(H)/H$ 可由一个\Term{separated}{分离}、\Term{finite}{有限}\Term{etale}{平展}的 $S$ 上\Term{scheme}{概形}表示。这由 5.3.18 以及 3.5 或第 XIX 讲 2.5 (ii) 得到。
\begin{SGAresult}{推论}{5.6.4}
设 $(G,T,M,R)$ 为一个\Term{split-group}{分裂}\Term{reductive-group}{约化群}。若 $R'\subset R$ 为\Term{closed-root-subset}{闭子集}且 $R'\cap -R'=\varnothing$，则 $R'$ 包含于一个\Term{positive-root-system}{正根系}中。
\end{SGAresult}
事实上，由 5.4.7，$R'$ 为 (R) 型，故结论由 5.6.1 得到。
\begin{SGAresult}{推论}{5.6.5}
在 5.6.1 的条件下，$G$ 中的乘法诱导\Term{isomorphism}{同构}
\[\prod_{\alpha\in R'}U_\alpha\xrightarrow{\sim}U_{R'},\]
其中 $U_{R'}$ 为 $G$ 的一个闭群子\Term{scheme}{概形}，在 $S$ 上\Term{smooth}{光滑}，\Term{geometric-fibre}{几何纤维}\Term{connected}{连通}且\Term{unipotent}{幺幂}，并且与 $R'$ 上次序的选择无关。此外，$H_{R'}$ 是\Term{semidirect-product}{半直积} $T\cdot U_{R'}$（$U_{R'}$ \Term{normal}{正规}）。
\end{SGAresult}
\begin{proof}
若 $R'\subset R_+$，则 $H_{R'}\cap U_{\ge1}$（5.5.6 的记号）是 $G$ 的一个\Term{finite-presentation}{有限呈示}\Term{closed-subgroup}{闭子群}，在 $H_{R'}$ 中\Term{normal}{正规}。由 5.6.1 (iv)，有 $H_{R'}=T\cdot U_{R'}$，其余断言由此得到。
\end{proof}
\textbf{注记 5.6.6。} 特别地，$U_{R_+}$ 就是 5.5.6 的群 $U_{\ge1}$。
现在从上述结果中提出几个关于 $U_{R'}$ 型群的推论。
\begin{SGAresult}{推论}{5.6.7}
设 $(G,T,M,R)$ 为一个\Term{split-group}{分裂}\Term{reductive-group}{约化群}，$R',R''$ 为 $R$ 的两个 (R) 型\Term{subset}{子集}，且 $R'\cap -R'=\varnothing$。

(i) 有 $U_{R'}\cap\operatorname{Norm}_G(H_{R''})=U_{R'\cap R''}$。

(ii) 设 $R'$ 为\Term{closed-root-subset}{闭子集}。若对所有满足 $\alpha+\beta\in R$ 的 $\alpha\in R'$、$\beta\in R''$ 均有 $\alpha+\beta\in R'$，则 $H_{R''}$ \Term{normalization}{正规化} $U_{R'}$。
\end{SGAresult}
\SourcePage{147}
\begin{proof}
(i) 立即由 5.6.5 和 5.6.1 (vi) 得到。为证 (ii)，由 5.4.4，只须说明 $T$ 及每个 $U_\beta$（$\beta\in R''$）\Term{normalization}{正规化} $U_{R'}$。\footnote{编者注：事实上，因 $G$ 的\Term{fiber}{纤维}\Term{connected}{连通}，由 VI$_B$ 5.5，它在 $S$ 上\Term{separated}{分离}；故由 XI 6.11（亦见 VI$_B$ 中新增的 6.5.5），$\operatorname{Norm}_G(U_{R'})$ 由 $G$ 的一个闭群子\Term{scheme}{概形} $N$ 表示。若 $N$ 包含 $T$ 及所有 $U_\beta$（$\beta\in R''$），则也包含\Term{big-cell}{大胞腔} $\Omega_{R_+,R''}$；由 5.4.4 和 EGA IV$_3$，11.10.10，该\Term{big-cell}{大胞腔}在 $H_{R''}$ 中\Term{schematically-dense}{概形稠密}（$H_{R''}$ 的\Term{fiber}{纤维}整，且 $\Omega_{R_+,R''}$ 含每条\Term{fiber}{纤维}的一个\Term{nonempty}{非空}\Term{open-subset}{开集}）。因此 $H_{R''}\subset N$，从而 $H_{R''}$ \Term{normalization}{正规化} $U_{R'}$。} 对 $T$ 显然；对 $U_\beta$，则由 5.5.2 和第 XXI 讲 3.1.2 得到。\footnote{编者注：所须证明的是，在 (ii) 的假设下，若 $\alpha\in R'$、$\beta\in R''$，则所有形如 $i\alpha+j\beta$（$i,j\in\mathbb N^*$）的根都属于 $R'$；为此，将对 XXI 2.3.5 的引用改为 XXI 3.1.2。也可直接考察秩为 $2$ 的\Term{root-system}{根系}得到此结论。}
\end{proof}
\begin{SGAresult}{推论}{5.6.8}
设 $(G,T,M,R)$ 为一个 $S$ 上的\Term{split-group}{分裂群}，$R_+$ 为一个\Term{positive-root-system}{正根系}，$\alpha$ 为 $R_+$ 的一个\Term{simple-root}{单根}（即使 $R_+-\{\alpha\}$ 为\Term{closed-root-subset}{闭子集}的 $R_+$ 中元素）。记
\[U_{\widehat\alpha}=U_{R_+-\{\alpha\}}.\]
则 (i) $U_{\widehat\alpha}$ 为 $B_{R_+}$ 的\Term{normal-subgroup}{正规子群}；(ii) $U_{R_+}$ 是 $U_{\widehat\alpha}$ 与 $U_\alpha$ 的\Term{semidirect-product}{半直积}；(iii) $U_{-\alpha}$ \Term{normalization}{正规化} $U_{\widehat\alpha}$；(iv) $Z_\alpha$ \Term{normalization}{正规化} $U_{\widehat\alpha}$。
若同样定义 $U_{-\widehat\alpha}=U_{R_--\{-\alpha\}}$（其中 $R_-=-R_+$），则
\[\Omega_{R_+}=U_{-\widehat\alpha}\cdot U_{-\alpha}\cdot T\cdot U_\alpha\cdot U_{\widehat\alpha}.\]
\end{SGAresult}
\begin{proof}
(ii) 由 5.6.5 得到，(i) 由 5.6.7 (ii) 得到。同样，(iii) 由 5.5.2 得到（事实上，若 $\beta\in R_+$、$\beta\ne\alpha$，则没有任何 $i(-\alpha)+j\beta$（$i,j>0$）会为负，因为 $\beta$ 至少含一个异于 $\alpha$ 的\Term{simple-root}{单根}）。因 $U_{-\alpha}\cdot T\cdot U_\alpha$ 在 $Z_\alpha$ 中\Term{schematically-dense}{概形稠密}，(iv) 由 (i)、(iii) 得到。最后一个断言由 (ii) 及其对 $U_{R_-}$ 的对应结论得到。
\end{proof}
回到一般情形。
\begin{SGAresult}{命题}{5.6.9}
设 $S$ 为\Term{scheme}{概形}，$G$ 为 $S$ 上的\Term{reductive-group}{约化群}，$H$ 为一个\Term{geometric-fibre}{几何纤维}\Term{solvable}{可解}的 (R) 型\Term{subgroup}{子群}。

(i) $D_S(H)=\operatorname{Hom}_{S\text{-gr.}}(H,\mathbb G_{m,S})$ 可由一个 $S$ 上的\Term{twisted-constant-group}{扭常值群}表示，其在 $s\in S$ 处的类型为 $\mathbb Z^{\operatorname{rgred}(G_s)}$。\Term{biduality}{双对偶}\Term{morphism}{态射}（第 VIII 讲第 1 节）
\[f:H\longrightarrow D_S(D_S(H))\]
\Term{smooth}{光滑}且\Term{surjectivity}{满射}。

(ii) $f$ 的核 $H^u$ 是 $H$ 中在 $S$ 上\Term{smooth}{光滑}、\Term{geometric-fibre}{几何纤维}\Term{connected}{连通}且\Term{unipotent}{幺幂}的最大闭\Term{normal-subgroup}{正规子群}\Term{scheme}{概形}。称之为 $H$ 的\Term{unipotent-part}{幺幂部分}，亦记 $H^u=\operatorname{rad}_u(H)$。
\SourcePage{148}
此时 $H^u$ 也是 $H$ 的\Term{commutator-sheaf}{换位子层}：从 $H$ 到一个对 (fppf) \Term{separated}{分离}的 $S$ 上\Term{abelian-group}{交换群}值\Term{presheaf}{预层}的任意群\Term{morphism}{态射}都在 $H^u$ 上为零，故经由 $H/H^u=D_S(D_S(H))$ \Term{resolution}{分解}。

(iii) 若 $T$ 为 $H$ 的一个\Term{maximal-torus}{极大环面}，则\Term{morphism}{态射} $T\to H$ 诱导\Term{isomorphism}{同构} $D_S(H)\xrightarrow\sim D_S(T)$ 及 $T\xrightarrow\sim D_S(D_S(H))$。此外，$H$ 与 $H^u$ 和 $T$ 的\Term{semidirect-product}{半直积}等同。

(iv) 在 5.6.1 的情形中，若 $H=H_{R'}$，则 $H^u=U_{R'}$。
\end{SGAresult}
\begin{proof}
本命题的断言对\Term{etale}{平展}\Term{topology}{拓扑}是\Term{local-scheme}{局部}的（第 X 讲 5.5）。故可归结到 5.6.1 的情形。已知 $H_{R'}$ 是 $U_{R'}$ 与 $T$ 的\Term{semidirect-product}{半直积}（5.6.5）。证明 $U_{R'}$ 是 $H_{R'}$ 的\Term{commutator-sheaf}{换位子层}：因 $H_{R'}/U_{R'}=T$ 交换，只须证明 (ii) 中的任意群\Term{morphism}{态射} $\varphi:H_{R'}\to V$ 在 $U_{R'}$ 上为零。只须证明它在每个 $U_\alpha$（$\alpha\in R'$）上为零。若 $t\in T(S')$、$X\in\mathbf W(\mathfrak g^\alpha)(S')$，则
\[1=\varphi\bigl(t\exp_\alpha(X)t^{-1}\exp_\alpha(-X)\bigr)
=\varphi\bigl(\exp_\alpha((\alpha(t)-1)X)\bigr).\]
因 $\alpha:T\to\mathbb G_{m,S}$ \Term{faithfully-flat}{忠实平坦}，立即推出 $\varphi$ 在 $U_\alpha^\times$ 上为零；但 $U_\alpha$ 的任意\Term{section}{截面}\Term{local-scheme}{局部}均为两个 $U_\alpha^\times$ 的\Term{section}{截面}之和。因此，对任意上述 $V$，有
\[\operatorname{Hom}_{S\text{-gr.}}(H,V)=\operatorname{Hom}_{S\text{-gr.}}(H/U_{R'},V).\]
对 $V=\mathbb G_{m,S}$ 应用此结果，立即得到 (i)、(iii)，进而得到 (iv) 及 (ii) 的第二个断言。现在只须证明 (ii) 的第一个断言；唯一非显然之处是，$H$ 中每个在 $S$ 上\Term{smooth}{光滑}、\Term{geometric-fibre}{几何纤维}\Term{connected}{连通}且\Term{unipotent}{幺幂}的闭\Term{normal-subgroup}{正规子群} $U$ 都是 $H^u$ 的\Term{subgroup}{子群}。先有：
\end{proof}
\begin{SGAresult}{引理}{5.6.10}
设 $G$ 为 $S$ 上的\Term{reductive-group}{约化群}，$T$ 为一个\Term{maximal-torus}{极大环面}，$U$ 为 $G$ 的一个群子\Term{scheme}{概形}，在 $S$ 上\Term{smooth}{光滑}、\Term{geometric-fibre}{几何纤维}\Term{unipotent}{幺幂}，且由 $T$ \Term{normalization}{正规化}。则 $U\cap T=e$。
\end{SGAresult}
\begin{proof}
因 $T=\operatorname{Centr}_G(T)$，有 $U\cap T=U^T$（$\operatorname{int}(T)$ 下的\Term{invariants}{不变量}）。应用第 XIX 讲 1.4，推出 $U\cap T$ 在 $S$ 上\Term{smooth}{光滑}；它在 $S$ 上也\Term{radiciel}{根性}：对每个 $s\in S$，$U(\overline s)\cap T(\overline s)$ 由同时\Term{unipotent}{幺幂}且\Term{semisimple}{半单}的元素组成。引理得证。
\end{proof}
\begin{proof}[继续证明 5.6.9 (ii)]
若 $U$ 是上述 $H$ 的一个\Term{normal-subgroup}{正规子群}，则\Term{semidirect-product}{半直积} $T\cdot U$ 是 $G$ 的一个\Term{geometric-fibre}{几何纤维}\Term{solvable}{可解}的 (R) 型\Term{subgroup}{子群}。因此可设它形如 $H_{R''}$，其中 $R''\subset R'$。只须证明 $U=U_{R''}$，从而归结到 $H=T\cdot U$ 的情形；但因商 $H/U$ 交换，$U$ 是 $H$ 的\Term{commutator-sheaf}{换位子层} $H^u$ 的子\Term{sheaf}{层}。证毕。
\end{proof}
注意，我们实际上已证明：
\begin{SGAresult}{命题}{5.6.11}
设 $S$ 为\Term{scheme}{概形}，$G$ 为 $S$ 上的\Term{reductive-group}{约化群}，$T$ 为 $G$ 的一个\Term{maximal-torus}{极大环面}。映射
\[H\longmapsto H^u,\qquad U\longmapsto T\cdot U\]
\SourcePage{149}
在 $G$ 中包含 $T$、\Term{geometric-fibre}{几何纤维}\Term{solvable}{可解}的 (R) 型\Term{subgroup}{子群} $H$ 的\Term{set}{集合}，与 $G$ 中在 $S$ 上\Term{smooth}{光滑}、由 $T$ \Term{normalization}{正规化}、\Term{geometric-fibre}{几何纤维}\Term{connected}{连通}且\Term{unipotent}{幺幂}的\Term{subgroup}{子群} $U$ 的\Term{set}{集合}之间，给出互逆的\Term{bijection}{双射}。\footnote{编者注：已删去 $U$ 为闭的假设，因为该假设自动成立。事实上，对这样的 $U$，由 5.6.10 有 $U\cap T=e$，故\Term{semidirect-product}{半直积} $H=T\cdot U$ 是 $G$ 的一个\Term{geometric-fibre}{几何纤维}\Term{solvable}{可解}的 (R) 型\Term{subgroup}{子群}（参见 5.2.1）。故由 5.6.3，$H$ 在 $G$ 中闭；因 $U$ 在 $H$ 中闭，它在 $G$ 中也闭。}
特别地，若 $(G,T,M,R)$ \Term{split-group}{分裂}，则群 $H_{R'}$ 与 $U_{R'}$ 相互对应。
\end{SGAresult}
\begin{SGAresult}{推论}{5.6.12}
设 $S$ 为\Term{scheme}{概形}，$(G,T,M,R)$ 为一个 $S$ 上的\Term{split-group}{分裂群}（分别地，再设 $R_+$ 为 $R$ 的一个\Term{positive-root-system}{正根系}，定义\Term{borel-subgroup}{Borel 子群} $B$）。
$G$ 的每个\Term{geometric-fibre}{几何纤维}\Term{connected}{连通}且\Term{unipotent}{幺幂}、由 $T$ \Term{normalization}{正规化}的\Term{smooth}{光滑}群子\Term{scheme}{概形}（分别地，$B^u$ 的每个由 $T$ \Term{normalization}{正规化}的\Term{smooth}{光滑}群子\Term{scheme}{概形}），在 $S$ 上\Term{local-scheme}{局部}均形如 $U_{R'}$，其中 $R'$ 是 $R$ 的一个包含于某个\Term{positive-root-system}{正根系}中的 (R) 型\Term{subset}{子集}（分别地，是 $R_+$ 的一个 (R) 型\Term{subset}{子集}）。
\end{SGAresult}
对\Term{lie-bracket}{括号}中的相应情形，只须注意由 Bible，\S 13.2，定理 1 (d)，所给群的\Term{geometric-fibre}{几何纤维}\Term{unipotent}{幺幂}且\Term{connected}{连通}。
命题 5.6.9 还有下述推论：
\begin{SGAresult}{推论}{5.6.13}
设 $S$ 为\Term{scheme}{概形}，$G$ 为 $S$ 上的\Term{reductive-group}{约化群}，$H$ 为一个\Term{geometric-fibre}{几何纤维}\Term{solvable}{可解}的 (R) 型\Term{subgroup}{子群}，$\operatorname{Tor}(H)$ 为 $H$ 的\Term{maximal-torus}{极大环面}\Term{functor}{函子}：
\[\operatorname{Tor}(H)(S')=\{H_{S'}\text{ 的极大环面}\}.\]
则 $\operatorname{Tor}(H)$ 可由一个\Term{affine}{仿射}\Term{smooth}{光滑}的 $S$ 上\Term{scheme}{概形}表示；对于\Term{action}{作用}律 $(h,T)\mapsto\operatorname{int}(h)T$，它是 $H^u$ 下的一个\Term{principal-homogeneous-space}{主齐性纤维丛}。\footnote{编者注：这重新证明并精确化了 XII 1.10（对 $G$ \Term{reduced}{约化}的情形）。}
\end{SGAresult}
\begin{proof}
若 $T,T'$ 是 $H_{S'}$ 的两个\Term{maximal-torus}{极大环面}，则存在唯一的\Term{section}{截面} $h\in H^u(S')$ 使 $\operatorname{int}(h)T=T'$。$h$ 的唯一性立即由 $\operatorname{Norm}_G(T)\cap H^u=e$ 得到（例如参见 5.6.1）；故只须对\Term{etale}{平展}\Term{topology}{拓扑}\Term{local-scheme}{局部}证明其存在性。由 5.2.6 和 5.1.2 (a)，可设 $T,T'$ 由 $H$ 的一个\Term{section}{截面}\Term{conjugacy}{共轭}；因由 5.6.9 (iii) 有 $H=H^u\cdot T$，即得所求结论。因此 $\operatorname{Tor}(H)$ 是 $H^u$ 下的一个\Term{principal-homogeneous-space}{主齐性层}，而 $H^u$ 在 $S$ 上\Term{affine}{仿射}且\Term{smooth}{光滑}，于是立即得到题述结论。\footnote{编者注：根据 SGA 1，VIII 2.1 及 EGA IV$_4$，17.7.1。}
\end{proof}
\SGAsection{5.7}{Bruhat 定理}
\textbf{回顾 5.7.1。} 设 $k$ 为\Term{algebraically-closed-field}{代数闭域}，$G$ 为 $k$ 上的\Term{reductive-group}{约化群}，$B$ 为 $G$ 的一个\Term{borel-subgroup}{Borel 子群}，$T$ 为 $B$ 的一个\Term{maximal-torus}{极大环面}，$N=\operatorname{Norm}_G(T)$。则
\[G(k)=B(k)N(k)B(k);\]
\SourcePage{150}
这就是 Bruhat 定理（Bible，\S 13.4，定理 3 的推论 1）；更精确地，沿用 3.6 的记号，\Term{set}{集合}
\[B(k)N_w(k)B(k)=B^u(k)N_w(k)B^u(k)\]
在 $w$ 遍历 $(N/T)(k)$ 时给出 $G(k)$ 的一个划分。若 $B'$ 是 $G$ 中包含 $T$ 的另一个\Term{borel-subgroup}{Borel 子群}，则\Term{set}{集合} $B'(k)N_w(k)B(k)$ 也给出 $G(k)$ 的一个划分。事实上，若 $y\in N(k)$ 满足 $\operatorname{int}(y)B=B'$，则
\[yB(k)N_w(k)B(k)=B'(k)N_{yw}(k)B(k).\]
\begin{SGAresult}{定义}{5.7.2}
设 $(G,T,M,R)$ 为一个 $S$ 上的\Term{split-group}{分裂群}，$R_-$ 为 $R$ 的一个\Term{positive-root-system}{正根系}，$B'=B_{R_-}$ 为它定义的\Term{borel-subgroup}{Borel 子群}。对 $w\in W$，记（参见 5.6.5）
\[R_-^w=R_-\cap w(R_-),\qquad B_w^{\prime u}=U_{R_-^w}=\prod_{\alpha\in R_-^w}U_\alpha.\]
若 $n_w\in\operatorname{Norm}_G(T)(S)$ 为 $w$ 的一个代表元（3.8），也可写成
\[B_w^{\prime u}=B^{\prime u}\cap\operatorname{int}(n_w)B^{\prime u}.\]
\end{SGAresult}
\begin{SGAresult}{引理}{5.7.3}
设 $(G,T,M,R)$ 为一个 $S$ 上的\Term{split-group}{分裂群}，$R_+$ 为 $R$ 的一个\Term{positive-root-system}{正根系}，$R_-=-R_+$，$B$（分别为 $B'$）为 $R_+$（分别为 $R_-$）定义的 $G$ 的\Term{borel-subgroup}{Borel 子群}。设 $w\in W$，$N_w$、$B_w^{\prime u}$ 为 $G$ 的相应子\Term{scheme}{概形}（3.8 和 5.7.2）。

(i) 由 $G$ 中乘法诱导的\Term{morphism}{态射} $B'\times_SN_w\times_SB\to G$ 的像\Term{sheaf}{层} $B'\cdot N_w\cdot B$，可由 $G$ 的一个子\Term{scheme}{概形}表示（事实上是\Term{open-subset}{开集} $n_w\Omega_{R_+}$ 的一个闭子\Term{scheme}{概形}）。

(ii) 由 $G$ 中乘法诱导的\Term{morphism}{态射}
\[B_w^{\prime u}\times_SN_w\times_SB^u\longrightarrow G\]
是像为上述子\Term{scheme}{概形}的\Term{immersion}{浸入}。
\end{SGAresult}
\begin{proof}
先说明 (ii) 的\Term{morphism}{态射}为\Term{immersion}{浸入}。由定义，$\operatorname{int}(n_w)^{-1}$ 诱导 $B_w^{\prime u}$ 到 $B^{\prime u}$ 的一个\Term{closed-immersion}{闭浸入}，故 $(u,b)\mapsto n_w^{-1}un_wb$ 诱导\Term{closed-immersion}{闭浸入}
\[B_w^{\prime u}\times_SB\longrightarrow\Omega_{R_+}.\]
这立即推出 (ii) 的\Term{morphism}{态射}诱导其定义域到\Term{open-subset}{开集} $n_w\Omega_{R_+}$ 的一个\Term{closed-immersion}{闭浸入}。为证 (i)，只须说明
\[B'(S)N_w(S)B(S)=B_w^{\prime u}(S)N_w(S)B^u(S).\]
若 $\alpha\in R$，有 $\operatorname{int}(n_w)U_\alpha(S)=U_{w(\alpha)}(S)$，所以若 $w^{-1}(\alpha)\in R_+$，则
\SourcePage{151}
\[\begin{split}
U_\alpha(S)N_w(S)B(S)&=U_\alpha(S)n_wT(S)B^u(S)\\
&=n_wU_{w^{-1}(\alpha)}(S)T(S)B^u(S)\\
&=n_wB(S)=N_w(S)B^u(S).
\end{split}\]
结合 $R_-^w$ 的定义，立即得到所求断言。
\end{proof}
\begin{SGAresult}{定理}{5.7.4}
设 $S$ 为\Term{scheme}{概形}，$(G,T,M,R)$ 为一个 $S$ 上的\Term{split-group}{分裂群}，$B$ 为\Term{positive-root-system}{正根系} $R_+$ 定义的\Term{borel-subgroup}{Borel 子群}，$B'$ 为 $R_-=-R_+$ 定义的\Term{borel-subgroup}{Borel 子群}。

(i)（Bruhat 定理）$w$ 遍历 $W$ 时，\Term{scheme}{概形} $B_w^{\prime u}\cdot N_w\cdot B$ 给出 $G$ 的\Term{underlying-set}{底集}合的一个划分。

(ii) 对每个 $w\in W$，取其在 $\operatorname{Norm}_G(T)(S)$ 中的代表元 $n_w$（3.8）；则\Term{open-subset}{开集} $n_w\Omega=n_wB^{\prime u}\cdot T\cdot B^u$ 在 $w$ 遍历 $W$ 时构成 $G$ 的一个\Term{covering}{覆盖}。
\end{SGAresult}
事实上，两个断言均可在\Term{geometric-fibre}{几何纤维}上验证，并由 5.7.1 和 5.7.3 得到结论。

\textbf{注记 5.7.5。} (i) 推出，若 $S$ 为域的谱，则 $G(S)$ 为所有 $B_w^{\prime u}(S)\cdot T(S)\cdot B^u(S)$ 的\Term{disjoint-union}{不交并}。对任意 $S$（即使是\Term{local-scheme}{局部}或 Artin 的），相应断言显然不成立。不过注意，(ii) 推出：若 $S$ \Term{local-scheme}{局部}，则 $G(S)$ 确是所有 $n_w\Omega(S)$ 的并。事实上：
\begin{SGAresult}{推论}{5.7.6}
设 $\Delta$ 为\Term{local-scheme}{局部}\Term{scheme}{概形} $S$ 上的\Term{split-group}{分裂群} $G$ 的一个\Term{simple-root}{单根}系。

(i) $G(S)$ 由 $T(S)$ 及 $U_\alpha(S)$（$\alpha\in\Delta\cup-\Delta$）生成。

(ii) 若 $G$ \Term{simply-connected}{单连通}（4.3.3），则 $G(S)$ 已由 $U_\alpha(S)$（$\alpha\in\Delta\cup-\Delta$）生成。
\end{SGAresult}
\begin{proof}
设 $H$ 为 $G(S)$ 中由 $U_\alpha(S)$（$\alpha\in\Delta\cup-\Delta$）生成的\Term{subgroup}{子群}。先注意 $H$ 包含每个 $s_\alpha$（$\alpha\in\Delta$）在 $\operatorname{Norm}_G(T)(S)$ 中的一个代表元（第 XX 讲 3.1），故包含每个 $w\in W$ 的一个代表元 $n_w$。
因每个 $\alpha\in R$ 均可写成 $w(\alpha_0)$，其中 $w\in W$、$\alpha_0\in\Delta$，所以
\[U_\alpha(S)=\operatorname{int}(n_w)U_{\alpha_0}(S)\subset H.\]
故由 $H$ 和 $T(S)$ 生成的\Term{subgroup}{子群}包含 $\Omega(S)$，由前述注记，它就是整个 $G(S)$。
现在若 $G$ \Term{simply-connected}{单连通}，证明 $H\supset T(S)$。由第 XX 讲 2.7，对每个 $\alpha\in\Delta$，$H$ 包含 $\alpha^*(\mathbb G_m(S))$；只须应用 4.3.8。
\end{proof}
\textbf{注记 5.7.6.1。} 对每个 $\alpha\in\Delta$，可以只取 $U_\alpha(S)$ 及对称 $s_\alpha$ 的一个代表元 $w_\alpha$，或取 $U_\alpha(S)$ 及 $U_{-\alpha}^\times$ 的一个\Term{section}{截面}，等等，来代替 $U_\alpha(S)$ 和 $U_{-\alpha}(S)$。
\begin{SGAresult}{推论}{5.7.7}
若 $G$ 的\Term{semisimple-rank}{半单秩}为 $1$，取 $u_\alpha\in U_\alpha^\times(S)$。则 $\Omega$ 与 $u_\alpha\Omega$ 构成 $G$ 的一个\Term{covering}{覆盖}。
\end{SGAresult}
\SourcePage{152}
\begin{proof}
若 $u_{-\alpha}$ 为与 $u_\alpha$ \Term{pairing}{配对}的 $U_{-\alpha}$ 的\Term{section}{截面}（参见 1.3），则由 5.7.4 (ii)，有
\[G=\Omega\cup u_{-\alpha}^{-1}u_\alpha u_{-\alpha}^{-1}\Omega,\]
从而 $G=u_{-\alpha}G=u_{-\alpha}\Omega\cup u_\alpha u_{-\alpha}^{-1}\Omega=\Omega\cup u_\alpha\Omega$。
\end{proof}
\begin{SGAresult}{推论}{5.7.8}
设 $S$ 为\Term{scheme}{概形}，$G$ 为 $S$ 上的\Term{reductive-group}{约化群}。则 $G$ 在 $S$ 上\Term{essentially-free}{本质自由}（第 VIII 讲 6.1）。\footnote{编者注：亦见 VI$_B$ 中新增的 6.2.1 至 6.2.6 及 6.5.2 至 6.5.5。}
\end{SGAresult}
\begin{proof}
该断言对 (fpqc) \Term{topology}{拓扑}是\Term{local-scheme}{局部}的，可设 $G$ \Term{split-group}{分裂}。此时 $G$ 具有由\Term{isomorphism}{同构}于 $\mathbb G_{a,S}^N\times_S\mathbb G_{m,S}^n$、因而\Term{essentially-free}{本质自由}的\Term{open-subset}{开集}构成的\Term{covering}{覆盖}。
\end{proof}
\begin{SGAresult}{引理}{5.7.9}
在 5.7.4 的条件下，设 $\alpha$ 为 $R_+$ 的一个\Term{simple-root}{单根}，$u_\alpha\in U_\alpha^\times(S)$。对每个 $v\in U_{-\alpha}(S)$，有 $\Omega\cdot v\subset\Omega\cup u_\alpha\cdot\Omega$。
\end{SGAresult}
\begin{proof}
这是比较 $G$ 的两个\Term{open-subset}{开集}，只须在 $S$ 是域 $k$ 的谱时证明。故须证明 $\Omega(k)v\subset\Omega(k)\cup u_\alpha\Omega(k)$。利用 5.6.8 的\Term{resolution}{分解}，有
\[\begin{split}
\Omega(k)v&=B^{\prime u}(k)T(k)B^u(k)v\\
&=U_{-\widehat\alpha}(k)U_{-\alpha}(k)T(k)U_\alpha(k)U_{\widehat\alpha}(k)v\\
&\subset U_{-\widehat\alpha}(k)Z_\alpha(k)U_{\widehat\alpha}(k)v.
\end{split}\]
应用 5.6.8 (iii)，再对群 $Z_\alpha$ 利用 5.7.7，得到
\[\begin{split}
\Omega(k)v&\subset U_{-\widehat\alpha}(k)Z_\alpha(k)vU_{\widehat\alpha}(k)
\subset U_{-\widehat\alpha}(k)Z_\alpha(k)U_{\widehat\alpha}(k)\\
&\subset U_{-\widehat\alpha}(k)U_{-\alpha}(k)T(k)U_\alpha(k)U_{\widehat\alpha}(k)\\
&\qquad\cup U_{-\widehat\alpha}(k)u_\alpha U_{-\alpha}(k)T(k)U_\alpha(k)U_{\widehat\alpha}(k).
\end{split}\]
再次利用 5.6.8 (iii)（将 $R_+$ 换为 $R_-$），即得结论。
\end{proof}
\begin{SGAresult}{命题}{5.7.10}
在 5.7.4 的条件下，对每个\Term{simple-root}{单根} $\alpha$ 取 $u_\alpha\in U_\alpha^\times(S)$。令 $U_1$ 为 $B^u(S)$ 中由所有 $u_\alpha$ 生成的子\Term{monoid}{幺半群}。\Term{open-subset}{开集} $u\Omega$（$u\in U_1$）构成 $G$ 的一个\Term{covering}{覆盖}。
\end{SGAresult}
\begin{proof}
仍可设 $S$ 是域 $k$ 的谱；由 5.7.6，只须证明 $\bigcup_{u\in U_1}u\Omega(k)$ 在右乘 $T(k)$、$U_\alpha(k)$、$U_{-\alpha}(k)$（$\alpha$ 为\Term{simple-root}{单根}）时\Term{stable}{稳定}。前两种情形显然，最后一种由引理得到。
\end{proof}
\SourcePage{153}
\textbf{注记 5.7.11。} 指出 5.7.2 的一个特例。若 $w=s_\alpha$ 为关于\Term{simple-root}{单根} $\alpha$ 的对称，则 $R_-\cap s_\alpha(R_-)=R_--\{-\alpha\}$（第 XXI 讲 3.3.1），故沿用 5.6.8 的记号，有 $B_{s_\alpha}^{\prime u}=U_{-\widehat\alpha}$。

\textbf{注记 5.7.12。} 实际上，5.7.10 的证明立即给出：在 5.7.10 的条件下，设 $\Gamma$ 为 $G(S)$ 的一个子\Term{monoid}{幺半群}；\Term{open-subset}{开集} $g\Omega$（$g\in\Gamma$）构成 $G$ 的一个\Term{covering}{覆盖}，当且仅当对每个 $s\in S$ 和每个\Term{simple-root}{单根} $\alpha$，有
\[(u_\alpha)_sB^{\prime u}(\overline s)\subset\Gamma\cdot B^{\prime u}(\overline s)\cdot T(\overline s)\cdot B^u(\overline s).\]
\textbf{注记 5.7.13。} 由 5.5.5 (iii)，按 5.7.1 的方式论证，立即得到 5.7.4 的下述变式：设 $(G,T,M,R)$ 为一个 $S$ 上的\Term{split-group}{分裂群}，$B,B'$ 为 $G$ 中包含 $T$ 的两个\Term{borel-subgroup}{Borel 子群}；对每个 $w\in W$，\Term{sheaf}{层} $B'\cdot N_w\cdot B$ 可由 $G$ 的一个子\Term{scheme}{概形}表示；这些子\Term{scheme}{概形}在 $w\in W$ 时构成 $G$ 的\Term{underlying-set}{底集}合的一个划分。也可给出 5.7.3 (ii) 的对应\Term{form}{形式}：只须置
\[B_w^{\prime u}=B^{\prime u}\cap\operatorname{int}(n_w)\widetilde B^u,\]
其中 $\widetilde B$ 是关于 $T$ 与 $B$ “\Term{opposite-borel}{相反}”的\Term{borel-subgroup}{Borel 子群}（参见 5.9.2）。
\begin{SGAresult}{命题}{5.7.14}
设 $S$ 为\Term{scheme}{概形}，$G$ 为 $S$ 上的\Term{reductive-group}{约化群}，$\operatorname{Ad}:G\to\operatorname{GL}_{\mathcal O_S}(\mathfrak g)$ 为其\Term{adjoint-representation}{伴随表示}。则 $\operatorname{Ker}(\operatorname{Ad})=\operatorname{Centr}(G)$；换言之，由 $\operatorname{Ad}$ 取商得到的\Term{canonical-homomorphism}{典范同态}
\[\overline{\operatorname{Ad}}:G/\operatorname{Centr}(G)=\operatorname{ad}(G)\longrightarrow\operatorname{GL}_{\mathcal O_S}(\mathfrak g)\]
是\Term{monomorphism}{单态射}。\footnote{编者注：由第 XVI 讲 1.5 (a)，它甚至是\Term{closed-immersion}{闭浸入}。}
\end{SGAresult}
\begin{proof}
可设 $G$ \Term{split-group}{分裂}。在 $\Gamma_0(R)$ 上取一个与\Term{group-structure}{群结构}相容的全序，令 $R_+$ 为\Term{positive-root}{正根}\Term{set}{集合}。由 5.7.4 (ii) 和 4.1.6，只须证明：若 $n_w$ 为 $w\in W$ 的一个代表元，$u\in U(S)$、$t\in T(S)$、$v\in U^-(S)$，且 $\operatorname{Ad}(n_wvtu)=\mathrm{id}$，则 $w=e$、$v=e$、$u=e$。对每个 $m\in R\cup\{0\}$，置
\[\mathfrak g^{>m}=\bigoplus_{n>m}\mathfrak g^n,\qquad\mathfrak g^{<m}=\bigoplus_{n<m}\mathfrak g^n.\]
设 $X\in\Gamma(S,\mathfrak g^m)$；写成 $\operatorname{Ad}(tu)X=\operatorname{Ad}(v^{-1}n_w^{-1})X$。有
\[\operatorname{Ad}(t)\operatorname{Ad}(u)X-m(t)X\in\Gamma(S,\mathfrak g^{>m}),\]
\[\operatorname{Ad}(v^{-1}n_w^{-1})X-\operatorname{Ad}(n_w^{-1})X\in\Gamma(S,\mathfrak g^{<w^{-1}(m)}).\]
\SourcePage{154}
若 $w\ne e$，存在 $\alpha\in R$ 使 $w^{-1}(\alpha)<\alpha$，令 $m=\alpha$，便得矛盾，因为
\[\operatorname{Ad}(tu)X\in\Gamma(S,\mathfrak g^\alpha+\mathfrak g^{>\alpha})\cap\Gamma(S,\mathfrak g^{w^{-1}(\alpha)}+\mathfrak g^{<w^{-1}(\alpha)})=0.\]
所以 $w=e$，可取 $n_w=e$；于是
\[\operatorname{Ad}(v^{-1})X-X\in\Gamma(S,\mathfrak g^{<m}\cap(\mathfrak g^m+\mathfrak g^{>m}))=0.\]
从而对每个 $X\in\Gamma(S,\mathfrak g^m)$，均有 $\operatorname{Ad}(v)X=X$，故 $\operatorname{Ad}(v)=\mathrm{id}$。同样 $\operatorname{Ad}(u)=\mathrm{id}$。再由 5.6.2 bis 得到结论。
\end{proof}
\SGAsection{5.8}{与约化群相伴的概形}
\begin{SGAresult}{定理}{5.8.1}
设 $S$ 为\Term{scheme}{概形}，$G$ 为 $S$ 上的\Term{reductive-group}{约化群}。令 $\mathcal H$ 为 $G$ 的 (R) 型\Term{subgroup}{子群}\Term{functor}{函子}：对每个 $S'\to S$，$\mathcal H(S')$ 为 $G_{S'}$ 的 (R) 型\Term{subgroup}{子群}的\Term{set}{集合}（参见 5.2.1）。则 $\mathcal H$ 可由一个在 $S$ 上\Term{quasi-projective}{拟射影}且\Term{finite-presentation}{有限呈示}的 $S$ 上\Term{scheme}{概形}表示。
\end{SGAresult}
\begin{proof}
\footnote{编者注：下文对原版作了详细补充。}令 $G'=G/\operatorname{Centr}(G)$ 为 $G$ 的\Term{adjoint-group}{伴随群}（4.3.6），$u$ 为\Term{morphism}{态射} $G\to G'$。由第 XII 讲 7.12，映射 $H'\mapsto u^{-1}(H')$ 在 $G'$ 与 $G$ 的 (R) 型\Term{subgroup}{子群}的\Term{set}{集合}之间给出\Term{bijection}{双射}，且这在任意\Term{base-change}{基变换}后仍成立。因此，以 $G'$ 代替 $G$，可设 $G$ 为\Term{adjoint-group}{伴随群}。考虑\Term{morphism}{态射}
\[u:\mathcal H\longrightarrow\operatorname{Grass}(\mathfrak g),\]
将每个 (R) 型\Term{subgroup}{子群}映到其\Term{lie-algebra}{Lie 代数}（它是 $\mathfrak g$ 的一个\Term{local-scheme}{局部}\Term{direct-summand}{直和因子}\Term{submodule}{子模}\footnote{编者注：参见第 II 讲 4.11.8。}）。由 5.3.3，$u$ 为\Term{monomorphism}{单态射}。只须证明它可由\Term{finite-presentation}{有限呈示}\Term{immersion}{浸入}表示，亦即证明：对每个 $S_1\to S$ 及给定的 $\mathfrak g_{S_1}$ 的一个\Term{local-scheme}{局部}\Term{direct-summand}{直和因子}\Term{submodule}{子模} $\mathfrak h$，使 $\mathfrak h_{S'}$ 成为 $G_{S'}$ 的某个 (R) 型\Term{subgroup}{子群}的\Term{lie-algebra}{Lie 代数}的那些 $S'\to S_1$，恰为经由 $S_1$ 的某个\Term{finite-presentation}{有限呈示}子\Term{scheme}{概形} $\Sigma$ \Term{resolution}{分解}的\Term{morphism}{态射}。用 $S$ 代替 $S_1$，归结到 $S_1=S$；还可设 $S$ \Term{affine}{仿射}。于是存在一个\Term{affine}{仿射} Noether \Term{scheme}{概形} $S_0$，使 $G$（分别为 $\mathfrak h$）由 $S_0$ 上的一个\Term{adjoint-group}{伴随}\Term{reductive-group}{约化群} $G_0$（分别为 $\mathfrak g_0=\operatorname{Lie}(G_0/S_0)$ 的一个\Term{local-scheme}{局部}\Term{direct-summand}{直和因子}\Term{submodule}{子模} $\mathfrak h_0$）经\Term{base-change}{基变换}得到。只须证明 $S_0$ 中存在具所求性质的子\Term{scheme}{概形} $\Sigma_0$（此时取 $\Sigma=\Sigma_0\times_{S_0}S$）。以 $S_0$ 代替 $S$，故可设 $S$ \Term{affine}{仿射}且 Noether（注意此时 $S$ 的每个子\Term{scheme}{概形}在 $S$ 上均\Term{finite-presentation}{有限呈示}）。最后，将 $S$ 换为足够小的\Term{open-subset}{开集}，可设 $\mathfrak g$ 自由秩为 $n$，而 $\mathfrak h$ 为自由秩 $r$ 的\Term{direct-summand}{直和因子}。
首先须表达 $\mathfrak h_{S'}$ 是 $\mathfrak g_{S'}$ 的 \Term{lie-subalgebra}{Lie 子代数}，即由 Lie \Term{lie-bracket}{括号}诱导的\Term{morphism}{态射}
\[\varphi:\mathfrak h\otimes\mathfrak h\xrightarrow{[\ ,\ ]}\mathfrak g\]
\SourcePage{155}
经由 $\mathfrak h$ \Term{resolution}{分解}。若 $(e_1,\ldots,e_n)$ 为 $\mathfrak g$ 的一组基，且 $(e_1,\ldots,e_r)$ 为 $\mathfrak h$ 的一组基，则 $\varphi$ 由 $\mathcal O_S$ 的\Term{section}{截面} $a_k^{ij}$（$i,j=1,\ldots,r$，$k=1,\ldots,n$）给出；上述条件\Term{equivalence}{等价}于 $S'\to S$ 经由方程 $a_k^{ij}=0$（$k=r+1,\ldots,n$，$i,j=1,\ldots,r$）定义的 $S$ 的闭子\Term{scheme}{概形}\Term{resolution}{分解}。以此闭子\Term{scheme}{概形}代替 $S$。
于是由 5.3.0，$N=\operatorname{Norm}_G(\mathfrak h)$ 是 $G$ 的一个闭群子\Term{scheme}{概形}，在 $G$ 上\Term{finite-presentation}{有限呈示}。现在须表达（参见 5.3.1）：$N_{S'}$ 在单位\Term{section}{截面}的每一点\Term{smooth}{光滑}，\Term{relative-dimension}{相对维数}为 $r=\operatorname{rang}(\mathfrak h)$，且包含 $\mathfrak h_{S'}\subset\operatorname{Lie}(N_{S'}/S')$ 为等式。
\footnote{编者注：原版接着指出，若用 $\mathfrak n_{N/G}$ 表示余法\Term{sheaf}{层} $\mathcal N_{N/G}$ 在单位\Term{section}{截面} $S\to N$ 下的\Term{inverse-image}{逆像}，则 $\mathfrak n_{N_{S'}/G_{S'}}\to\omega^1_{G_{S'}/S'}$ \Term{universally-injective}{普遍单射}这一条件，\Term{equivalence}{等价}于 $S'\to S$ 经由 $S$ 的某个开子\Term{scheme}{概形}\Term{resolution}{分解}。我们未能说明这一点，因为 $\mathcal N_{N/G}$ 的形成与\Term{base-change}{基变换}不交换；故用下面由 O. Gabber 提出的论证代替。}
因 $N$ 在 $S$ 上\Term{affine}{仿射}（它在 $G$ 中闭），单位\Term{section}{截面} $\varepsilon:S\to N$ 是\Term{closed-immersion}{闭浸入}，故 $\varepsilon(S)$ 由 $\mathcal A(N)$ 的一个\Term{quasi-coherent}{拟凝聚}\Term{ideal}{理想} $\mathcal J$ 定义。注意 $\mathfrak n_{S/N}=\mathcal J/\mathcal J^2$ 与 $\varepsilon^*(\Omega^1_{N/S})$ 等同，所以其形成与任意\Term{base-change}{基变换} $S'\to S$ 交换。由 EGA IV$_4$，17.12.1 中的\Term{equivalence}{等价}性 (c')$\Leftrightarrow$(a)（应用于 $f:N\to S$ 和 $j=\varepsilon$），$N_{S'}\to S'$ 在 $\varepsilon(S')$ 的每一点\Term{smooth}{光滑}且\Term{relative-dimension}{相对维数}为 $r$，当且仅当 $\mathfrak n_{S'/N'}=\mathfrak n_{S/N}\otimes_{\mathcal O_S}\mathcal O_{S'}$ \Term{locally-free}{局部自由}秩为 $r$，且对每个 $n\ge1$，\Term{morphism}{态射}
\[\varphi_n(S'):\operatorname{Sym}^n(\mathfrak n_{S'/N'})\longrightarrow\mathcal J_{S'}^n/\mathcal J_{S'}^{n+1}\]
为\Term{isomorphism}{同构}。记 $K_n(S')=\operatorname{Ker}\varphi_n(S')$。由 TDTE I，引理 3.6，$\mathfrak n_{S/N}\otimes_{\mathcal O_S}\mathcal O_{S'}$ \Term{locally-free}{局部自由}秩为 $r$，当且仅当 $S'\to S$ 经由 $S$ 的某个子\Term{scheme}{概形} $Z$ \Term{resolution}{分解}。以 $Z$ 代替 $S$，可设 $\mathfrak n_{S/N}=\mathcal J/\mathcal J^2$ \Term{locally-free}{局部自由}秩为 $r$。于是对每个 $S'\to S$，有
\[(\mathcal J^2/\mathcal J^3)\otimes_{\mathcal O_S}\mathcal O_{S'}=\mathcal J_{S'}^2/\mathcal J_{S'}^3,\qquad K_2(S')=K_2(S)\otimes_{\mathcal O_S}\mathcal O_{S'}.\]
因此，$\varphi_2(S')$ 是\Term{isomorphism}{同构}，当且仅当 $S'\to S$ 经由 $S$ 的闭子\Term{scheme}{概形} $S_2$ \Term{resolution}{分解}，其中 $S_2$ 由 $\operatorname{Sym}^2(\mathfrak n_{S/N})^*\otimes K_2(S)$ 在 $\mathcal O_S$ 中的像所生成的\Term{ideal}{理想}定义。在 $S_2$ 上，$\mathcal J^2/\mathcal J^3$ 与 $\operatorname{Sym}^2(\mathfrak n_{S/N})$ \Term{isomorphism}{同构}，故\Term{locally-free}{局部自由}；同样的论证说明，$\varphi_3(S')$ 是\Term{isomorphism}{同构}，当且仅当 $S'\to S$ 经由 $S$ 的某个闭子\Term{scheme}{概形} $S_3$ \Term{resolution}{分解}，依此类推。于是 $N_{S'}\to S'$ 在 $\varepsilon(S')$ 的每一点\Term{smooth}{光滑}且\Term{relative-dimension}{相对维数}为 $r$，当且仅当 $S'\to S$ 经由所有 $S_n$ 的交这个闭子\Term{scheme}{概形} $Z$ \Term{resolution}{分解}。但此时，对每个 $S'\to Z$，$\operatorname{Lie}(N_{S'}/S')$ \Term{local-scheme}{局部}是 $\mathfrak g_{S'}$ 的秩 $r$ 的\Term{direct-summand}{直和因子}，故包含 $\mathfrak h_{S'}\subset\operatorname{Lie}(N_{S'}/S')$ 是等式。于是置 $H=N^0$。
以 $Z$ 代替 $S$，只须再表达：对每个 $s'\in S'$，$H_{s'}$ 与 $G_{s'}$ 的\Term{reductive-rank}{约化秩}相同；\Term{equivalence}{等价}地，对 $S'$ 在 $S$ 中的\Term{set}{集合}意义的像中的每个点 $s$，$H_s$ 与 $G_s$ 的\Term{reductive-rank}{约化秩}相同。而该条件定义 $S$ 的一个\Term{open-subset}{开子集}（第 XIX 讲 6.2）。
\end{proof}
\textbf{注记。} 一般而言，\Term{scheme}{概形} $\mathcal H$ 在 $S$ 上不\Term{smooth}{光滑}。但若 $6$ 在 $S$ 上\Term{invertible}{可逆}，或存在素数 $p$ 使 $p\cdot1_S=0$（即 $S$ 的\Term{characteristic}{特征}为 $p>0$），则它\Term{smooth}{光滑}。
\SourcePage{156}
\begin{SGAresult}{推论}{5.8.2}
设 $S$ 为\Term{scheme}{概形}，$G$ 为 $S$ 上的\Term{reductive-group}{约化群}，$H$ 为 $G$ 的一个 (R) 型\Term{subgroup}{子群}。（回顾 5.3.10：$\operatorname{Norm}_G(H)$ 可由 $G$ 的一个在 $S$ 上\Term{smooth}{光滑}的闭群子\Term{scheme}{概形}表示。）
则商\Term{sheaf}{层} $G/\operatorname{Norm}_G(H)$ 可由一个在 $S$ 上\Term{quasi-projective}{拟射影}、\Term{smooth}{光滑}且\Term{finite-presentation}{有限呈示}的 $S$ 上\Term{scheme}{概形}表示（事实上是 $\mathcal H$ 的一个\Term{open-subset}{开集}）。
\end{SGAresult}
\begin{proof}
考虑逐点由 $f(g)=\operatorname{int}(g)H$ 定义的\Term{morphism}{态射} $f:G\to\mathcal H$。由 5.3.9，此\Term{morphism}{态射}\Term{smooth}{光滑}且\Term{finite-presentation}{有限呈示}，故为开映射。令 $V=f(G)$，赋予其作为 $\mathcal H$ 的开子\Term{scheme}{概形}的结构。\Term{morphism}{态射} $f:G\to V$ 为\Term{covering}{覆盖}，核为 $\operatorname{Norm}_G(H)$；这说明 $G/\operatorname{Norm}_G(H)$ 由 $V$ 表示（参见第 IV 讲 4.6.5）。
\end{proof}
\begin{SGAresult}{推论}{5.8.3}
设 $S$ 为\Term{scheme}{概形}，$G$ 为 $S$ 上的\Term{reductive-group}{约化群}。考虑由下式定义的\Term{functor}{函子} $\operatorname{Tor}(G)$、$\operatorname{Bor}(G)$、$\operatorname{Kil}(G)$：
\[\begin{aligned}
\operatorname{Tor}(G)(S')&=\{G_{S'}\text{ 的极大环面}\},\\
\operatorname{Bor}(G)(S')&=\{G_{S'}\text{ 的 Borel 子群}\},\\
\operatorname{Kil}(G)(S')&=\{G_{S'}\text{ 的 Killing 对（参见 5.3.13）}\}.
\end{aligned}\]
(i) $\operatorname{Tor}(G)$、$\operatorname{Bor}(G)$、$\operatorname{Kil}(G)$ 可由在 $S$ 上\Term{smooth}{光滑}且\Term{finite-presentation}{有限呈示}、\Term{geometric-fibre}{几何纤维}整的 $S$ 上\Term{scheme}{概形}表示；它们在 $S$ 上分别\Term{affine}{仿射}、\Term{projective}{射影}、\Term{affine}{仿射}。

(ii) \Term{canonical}{典范}\Term{morphism}{态射} $\operatorname{Kil}(G)\to\operatorname{Tor}(G)$（分别为 $\operatorname{Kil}(G)\to\operatorname{Bor}(G)$）是\Term{finite}{有限}\Term{etale}{平展}\Term{surjectivity}{满射}（分别为\Term{affine}{仿射}\Term{smooth}{光滑}\Term{surjectivity}{满射}）。

(iii) 设 $T$ 为 $G$ 的一个\Term{maximal-torus}{极大环面}（分别地，$B$ 为 $G$ 的一个\Term{borel-subgroup}{Borel 子群}；分别地，$B\supset T$ 为 $G$ 的一个\Term{killing-pair}{Killing 对}）。由
\[g\mapsto\operatorname{int}(g)T,\qquad
 g\mapsto\operatorname{int}(g)B,\qquad
 g\mapsto(\operatorname{int}(g)B,\operatorname{int}(g)T)\]
定义的相应\Term{morphism}{态射}
\[G\to\operatorname{Tor}(G),\qquad G\to\operatorname{Bor}(G),\qquad G\to\operatorname{Kil}(G)\]
分别诱导\Term{isomorphism}{同构}
\[G/\operatorname{Norm}_G(T)\xrightarrow\sim\operatorname{Tor}(G),\quad
G/B\xrightarrow\sim\operatorname{Bor}(G),\quad
G/T\xrightarrow\sim\operatorname{Kil}(G).\]
\end{SGAresult}
\begin{proof}
首先，(iii) 由\Term{maximal-torus}{极大环面}（分别为\Term{borel-subgroup}{Borel 子群}、\Term{killing-pair}{Killing 对}）的\Term{conjugacy}{共轭}定理及等式
\[\operatorname{Norm}_G(B)=B,\qquad\operatorname{Norm}_G(B)\cap\operatorname{Norm}_G(T)=T\]
得到；这些结果均已证明（5.1.2、5.3.12、5.3.14、5.6.1）。因此\Term{canonical}{典范}\Term{morphism}{态射} $\operatorname{Tor}(G)\to\mathcal H$、$\operatorname{Bor}(G)\to\mathcal H$ 对\Term{etale}{平展}\Term{topology}{拓扑}\Term{local-scheme}{局部}可由\Term{open-immersion}{开浸入}表示（分别由 5.8.2、5.1.2 以及 5.5.5），故由\Term{descent}{下降}，$\operatorname{Tor}(G)$ 和 $\operatorname{Bor}(G)$ 可由 $\mathcal H$ 的\Term{open-subset}{开集}表示。同样，$\operatorname{Kil}(G)$ 对\Term{etale}{平展}\Term{topology}{拓扑}\Term{local-scheme}{局部}可由底上的\Term{affine}{仿射}\Term{scheme}{概形}表示（第 IX 讲 2.3），
\SourcePage{157}
故由\Term{affine}{仿射}\Term{scheme}{概形}的\Term{descent}{下降}，它可由一个 $S$ 上的\Term{affine}{仿射}\Term{scheme}{概形}表示。\footnote{编者注：参见 SGA 1，VIII 2.1。}
(ii) 的断言立即由 5.5.5 (ii) 和 5.6.13 得到。因此 $\operatorname{Tor}(G)$ 在 $S$ 上\Term{affine}{仿射}（EGA II 6.7.1）。只须再证明 $\operatorname{Bor}(G)$ 在 $S$ 上\Term{projective}{射影}。已知它\Term{quasi-projective}{拟射影}，只须证明其\Term{proper}{固有}；\footnote{编者注：可设 $S$ \Term{affine}{仿射}，再由 (i) 中 $\operatorname{Bor}(G)$ 在 $S$ 上\Term{finite-presentation}{有限呈示}，归结到 $S$ 是 Noether 的情形；此时满足 EGA IV$_3$，15.7.10 的假设。}因其\Term{fiber}{纤维}\Term{connected}{连通}，由 EGA IV$_3$，15.7.10，只须在\Term{geometric-fibre}{几何纤维}上证明。若 $S$ 是\Term{algebraically-closed-field}{代数闭域}的谱，则由 (iii) 有 $\operatorname{Bor}(G)=G/B$，由 Bible，\S 6.4，定理 4（或 [Ch05]，\S 6.5，定理 5）即得结论。
\end{proof}
\textbf{注记 5.8.4。} 在 5.8.3 的条件下，设 $Q$ 为 $G$ 的一个\Term{multiplicative-type}{乘法型}\Term{central-subgroup}{中心子群}。显然的\Term{morphism}{态射}定义\Term{isomorphism}{同构}
\[\operatorname{Tor}(G)\simeq\operatorname{Tor}(G/Q),\quad
\operatorname{Bor}(G)\simeq\operatorname{Bor}(G/Q),\quad
\operatorname{Kil}(G)\simeq\operatorname{Kil}(G/Q).\]
\begin{SGAresult}{推论}{5.8.5}
设 $S$ 为\Term{scheme}{概形}，$G$ 为 $S$ 上的\Term{reductive-group}{约化群}，$P$ 为 $G$ 的一个在 $S$ 上\Term{smooth}{光滑}且\Term{finite-presentation}{有限呈示}的群子\Term{scheme}{概形}。下列条件\Term{equivalence}{等价}：

(i) 对每个 $s\in S$，$P_s$ 是 $G_s$ 的一个\Term{parabolic-subgroup}{抛物子群}（即商\Term{scheme}{概形} $G_s/P_s$ 在 $s$ 上\Term{proper}{固有}；又\Term{equivalence}{等价}于 $P_s$ 包含 $G_s$ 的一个\Term{borel-subgroup}{Borel 子群}，参见 Bible，\S 6.4，定理 4，或 [Ch05]，\S 6.5，定理 5）。

(ii) 商\Term{sheaf}{层} $G/P$ 可由一个在 $S$ 上\Term{smooth}{光滑}且\Term{projective}{射影}的 $S$ 上\Term{scheme}{概形}表示。

此外，在这些条件下，$P$ 在 $G$ 中闭，\Term{fiber}{纤维}\Term{connected}{连通}，且 $P=\operatorname{Norm}_G(P)$。
\end{SGAresult}
\begin{proof}
显然 (ii)$\Rightarrow$(i)。若 (i) 成立，则 $P(\overline s)=\operatorname{Norm}_{G(\overline s)}(P_s)$，且 $P_s$ \Term{connected}{连通}（第一点参见 Bible，\S 12.3，引理 4；\footnote{编者注：下文已作详细补充。}第二点由此得到，因为 $P^0=P_s^0$ 是 $G_s$ 的一个由 $P_s$ \Term{normalization}{正规化}的\Term{parabolic-subgroup}{抛物子群}，故 $P^0(\overline s)=P(\overline s)$，于是 $P^0=P_s$）；所以 $P$ 为 (R) 型，且等于 $\operatorname{Norm}_G(P)$，故在 $G$ 中闭。由 5.8.2，$G/P=G/\operatorname{Norm}_G(P)$ 可由一个 $S$ 上的\Term{quasi-projective}{拟射影}\Term{scheme}{概形}表示。其\Term{fiber}{纤维}\Term{connected}{连通}且\Term{proper}{固有}，故由 5.8.3 的论证，它\Term{projective}{射影}。
\end{proof}
\textbf{注记 5.8.6。} 5.8.1、5.8.2、5.8.5 对 $S$ 上的 (RA) 型群，或满足 5.1.8 的 $S$ 上 (RR) 型群也成立。\footnote{编者注：事实上，5.8.1 的证明只用到 5.3.3（它对 (RA) 型群成立）及 XIX 6.2，而由 XII 1.7 (b)，后者也对 (RR) 型群成立。}

\textbf{注记 5.8.7。} 通过 $G$ 的\Term{inner-automorphism}{内自同构}，有\Term{canonical}{典范}\Term{action}{作用}
\[G\to\operatorname{Aut}_S(\operatorname{Tor}(G)),\quad
G\to\operatorname{Aut}_S(\operatorname{Bor}(G)),\quad
G\to\operatorname{Aut}_S(\operatorname{Kil}(G));\]
\SourcePage{158}
在 5.8.3 (iii) 的情形中，它们与\Term{canonical}{典范}\Term{action}{作用}
\[G\to\operatorname{Aut}_S(G/\operatorname{Norm}_G(T)),\quad G\to\operatorname{Aut}_S(G/B),\quad G\to\operatorname{Aut}_S(G/T)\]
等同。特别地可得
\[\begin{split}
\operatorname{Ker}\bigl(G\to\operatorname{Aut}_S(\operatorname{Tor}(G))\bigr)
&=\operatorname{Ker}\bigl(G\to\operatorname{Aut}_S(\operatorname{Bor}(G))\bigr)\\
&=\operatorname{Ker}\bigl(G\to\operatorname{Aut}_S(\operatorname{Kil}(G))\bigr)
=\operatorname{Centr}(G).
\end{split}\]
事实上，$\operatorname{Centr}(G)$ 显然在这三个\Term{scheme}{概形}上均\Term{action}{作用}平凡。反之，$G\to\operatorname{Aut}_S(\operatorname{Kil}(G))$ 的核就是 4.1.7 意义下“$G$ 的所有\Term{maximal-torus}{极大环面}的交”，故等于 $\operatorname{Centr}(G)$（同上）。对 $\operatorname{Bor}(G)$，注意“$G$ 的所有\Term{borel-subgroup}{Borel 子群}的交”也是“其所有\Term{maximal-torus}{极大环面}的交”（见下一小节）。对 $\operatorname{Tor}(G)$，利用第 XII 讲 4.11。
\SGAsection{5.9}{Borel 子群特有的性质}
这些性质大多将在第 XXVI 讲推广到\Term{parabolic-subgroup}{抛物子群}。
\begin{SGAresult}{定义}{5.9.1}
设 $S$ 为\Term{scheme}{概形}，$G$ 为 $S$ 上的\Term{reductive-group}{约化群}，$B,B'$ 为 $G$ 的两个\Term{borel-subgroup}{Borel 子群}。若 $B\cap B'$ 是 $G$ 的一个\Term{torus}{环面}（它必为\Term{maximal-torus}{极大环面}），则称 $B,B'$ 处于\Term{general-position}{一般位置}（或称 $B'$ 相对于 $B$ 处于\Term{general-position}{一般位置}）。
若 $T$ 为 $G$ 中包含于 $B$ 和 $B'$ 的一个\Term{maximal-torus}{极大环面}，且 $B\cap B'=T$，则称 $B,B'$（关于 $T$）\Term{opposite-borel}{相反}。
\end{SGAresult}
\begin{SGAresult}{命题}{5.9.2}
设 $S$ 为\Term{scheme}{概形}，$G$ 为 $S$ 上的\Term{reductive-group}{约化群}，$B$ 为 $G$ 的一个\Term{borel-subgroup}{Borel 子群}，$T$ 为 $B$ 的一个\Term{maximal-torus}{极大环面}。存在唯一的 $G$ 的\Term{borel-subgroup}{Borel 子群} $B'$，关于 $T$ 与 $B$ \Term{opposite-borel}{相反}。
若 $(G,T,M,R)$ 是 $G$ 关于 $T$ 的一个\Term{splitting}{分裂}，且 $B=B_{R_+}$（5.5.1），则 $B'=B_{-R_+}$。
\end{SGAresult}
\begin{proof}
由\Term{faithfully-flat}{忠实平坦}\Term{descent}{下降}，只须在\Term{splitting}{分裂}情形、$B=B_{R_+}$ 时证明（5.5.5 (iv)）。此时 $B_{-R_+}$ 确与 $B$ \Term{opposite-borel}{相反}（4.1.2）；证明它是 $G$ 中包含 $T$ 且与 $B$ \Term{opposite-borel}{相反}的唯一\Term{borel-subgroup}{Borel 子群}。若 $B'$ 是 $G$ 中包含 $T$ 的一个\Term{borel-subgroup}{Borel 子群}，则在 $S$ 上\Term{local-scheme}{局部}它形如 $B_{R'_+}$，其中 $R'_+$ 为 $R$ 的第二个\Term{positive-root-system}{正根系}（5.5.5 (iii)）。若 $R'_+\ne-R_+$，则存在 $\alpha\in R'_+\cap R_+$，因而 $U_\alpha\subset B_{R_+}\cap B_{R'_+}$。
\end{proof}
\begin{SGAresult}{命题}{5.9.3}
设 $S$ 为\Term{scheme}{概形}，$G$ 为 $S$ 上的\Term{reductive-group}{约化群}，$B$ 为 $G$ 的一个\Term{borel-subgroup}{Borel 子群}。

(i) 若 $B'$ 为 $G$ 的一个\Term{borel-subgroup}{Borel 子群}，则下列条件\Term{equivalence}{等价}：

(a) $B'$ 相对于 $B$ 处于\Term{general-position}{一般位置}（5.9.1）。

(b) $B^{\prime u}\cap B^u=e$。

(b') $B^{\prime u}\cap B=e$。

(c) $G$ 中乘法诱导一个\Term{open-immersion}{开浸入} $B^{\prime u}\times_SB\to G$。
\SourcePage{159}
(c') \Term{canonical}{典范}\Term{morphism}{态射} $B^{\prime u}\to G/B$ 是\Term{open-immersion}{开浸入}。

(ii) \Term{functor}{函子}
\[\operatorname{Opp}(B):S'\longmapsto
\{G_{S'}\text{ 中相对于 }B_{S'}\text{ 处于一般位置的 Borel 子群}\}\]
可由 $\operatorname{Bor}(G)$ 的一个开子\Term{scheme}{概形}表示（5.8.3）。由 $B'\mapsto B\cap B'$ 定义的\Term{morphism}{态射} $\operatorname{Opp}(B)\to\operatorname{Tor}(B)$ 是\Term{isomorphism}{同构}。特别地，由 5.6.13，$B^u$ 的\Term{inner-automorphism}{内自同构}赋予 $\operatorname{Opp}(B)$ 一个 $B^u$ 下的\Term{principal-homogeneous-space}{主齐性纤维丛}结构。
\end{SGAresult}
\begin{proof}
先考察 (i)。(a)$\Rightarrow$(c)：事实上 (c) 对\Term{etale}{平展}\Term{topology}{拓扑}是\Term{local-scheme}{局部}的，由 5.5.5 (iv)，归结到 $G$ 关于 $B\cap B'$ \Term{split-group}{分裂}且 $B$ 形如 $B_{R_+}$ 的情形；由 5.9.2，此时 $B^{\prime u}=U_{-R_+}$，归结到 4.1.2。
显然有 (c')$\Leftrightarrow$(c)$\Rightarrow$(b')$\Rightarrow$(b)。故只须证明 (b)$\Rightarrow$(a)。先证明 (ii)：第二个断言是 5.9.2 的\Term{form}{形式}推论，第三个由第二个及 5.6.13 立即得到；再证明第一个断言。它对\Term{etale}{平展}\Term{topology}{拓扑}是\Term{local-scheme}{局部}的，故可设 $B$ 具有一个\Term{maximal-torus}{极大环面} $T$；令 $B'_0$ 为关于 $T$ 与 $B$ \Term{opposite-borel}{相反}的\Term{subgroup}{子群}（5.9.2）。
由上述论证，\Term{canonical}{典范}\Term{morphism}{态射} $G\to G/B'_0\to\operatorname{Bor}(G)$（5.8.3）所诱导的 $B^u\to\operatorname{Bor}(G)$，诱导\Term{isomorphism}{同构} $B^u\xrightarrow\sim\operatorname{Opp}(B)$。所以有\Term{commutative-diagram}{交换图}
\[\begin{tikzcd}
B^u\arrow[r]\arrow[d,"\sim"']&G/B'_0\arrow[d,"\sim"]\\
\operatorname{Opp}(B)\arrow[r]&\operatorname{Bor}(G).
\end{tikzcd}\]
而 $B^u\to G/B'_0$ 是\Term{open-immersion}{开浸入}（由 (i) 的 (a)$\Rightarrow$(c')），故 (ii) 得证。立即记下其推论：
\end{proof}
\begin{SGAresult}{推论}{5.9.4}
设 $G$ 为 $S$ 上的\Term{reductive-group}{约化群}，$B,B'$ 为 $G$ 的两个\Term{borel-subgroup}{Borel 子群}。若 $s\in S$ 使 $B_s,B'_s$ 处于\Term{general-position}{一般位置}，则存在包含 $s$ 的 $S$ 的一个\Term{open-subset}{开集} $V$，使 $B_V,B'_V$ 处于\Term{general-position}{一般位置}。
\end{SGAresult}
\begin{proof}[继续证明 5.9.3]
现在只须证明 (b)$\Rightarrow$(a)。由上面的推论，只须在 $S$ 是\Term{algebraically-closed-field}{代数闭域} $k$ 的谱时证明。可设 $G$ 关于 $B$ 的一个\Term{maximal-torus}{极大环面} $T$ \Term{split-group}{分裂}。令 $B'_0$ 为与 $B$ \Term{opposite-borel}{相反}的\Term{borel-subgroup}{Borel 子群}。因 $G$ 的\Term{borel-subgroup}{Borel 子群}在 $G(k)$ 下\Term{conjugacy}{共轭}，存在 $g\in G(k)$ 使 $\operatorname{int}(g)B'_0=B'$。由 Bruhat 定理（5.7.4），可写成 $g=bnb'$，其中 $b\in B(k)$、$b'\in B'_0(k)$、$n\in\operatorname{Norm}_G(T)(k)$。故
\[B'=\operatorname{int}(b)\operatorname{int}(n)B'_0,\qquad
B'\cap B=\operatorname{int}(b)\bigl(\operatorname{int}(n)B'_0\cap B\bigr).\]
若 $n\notin T(k)$，则 $\operatorname{int}(n)B_0^{\prime u}\cap B^u\ne e$（参见 5.9.2 的证明）；所以 (b) 推出 $B'\cap B=\operatorname{int}(b)(B'_0\cap B)=\operatorname{int}(b)T$。证毕。
\end{proof}
\SourcePage{160}
\begin{SGAresult}{命题}{5.9.5}
设 $S$ 为\Term{scheme}{概形}，$G$ 为 $S$ 上的\Term{reductive-group}{约化群}，$B$ 为 $G$ 的一个\Term{borel-subgroup}{Borel 子群}，$B^u$ 为其\Term{unipotent-part}{幺幂部分}。存在 $B$ 的一列\Term{subgroup}{子群}
\[U_0=B^u\supset U_1\supset\cdots\supset U_n\supset\cdots\]
具有下列性质：

(i) 每个 $U_i$ 均\Term{smooth}{光滑}，\Term{fiber}{纤维}\Term{connected}{连通}，且在 $B$ 中为\Term{characteristic-subgroup}{特征子群}；$B^u$ 的\Term{inner-automorphism}{内自同构}在商（\Term{sheaf}{层}）$U_i/U_{i+1}$ 上\Term{action}{作用}平凡。

(ii) 对每个 $i\ge0$，存在\Term{locally-free}{局部自由} $\mathcal O_S$ 模 $\mathcal E_i$ 及 $S$ 上群\Term{sheaf}{层}的\Term{isomorphism}{同构} $U_i/U_{i+1}\xrightarrow\sim\mathbf W(\mathcal E_i)$。

(iii) 对每个 $s\in S$，当 $n\ge\dim(B_s^u)$ 时，$(U_n)_s=e$。
\end{SGAresult}
\begin{proof}
先设存在 $G$ 的一个\Term{splitting}{分裂} $(G,T,M,R)$ 及 $R$ 的一个\Term{positive-root-system}{正根系} $R_+$，使 $B=B_{R_+}$。记 $\Delta$ 为 $R_+$ 的\Term{simple-root}{单根}\Term{set}{集合}；对每个 $\alpha\in R_+$，用 $\operatorname{ord}(\alpha)$ 表示 $\alpha$ 在 $\Gamma_0(R)$ 的基 $\Delta$ 下各系数之和，称之为 $\alpha$ 关于 $R_+$ 的阶。有 $\operatorname{ord}(\alpha)\le\operatorname{Card}(R_+)$。对每个 $i>0$，令 $R^{(i)}$ 为阶大于 $i$ 的\Term{root}{根}的\Term{set}{集合}，这是一个\Term{positive-root}{正根}\Term{closed-root-subset}{闭集}，故\Term{constructible}{可构造}（5.6.5）$U_i=U_{R^{(i)}}$。
若 $\alpha\in R_+$ 且 $\beta\in R^{(i)}$，则 $\alpha+\beta\in R^{(i+1)}$。由 5.5.2，每个 $U_i$ 是 $B$ 的\Term{normal-subgroup}{正规子群}，而 $B^u$ 的\Term{inner-automorphism}{内自同构}在 $U_i/U_{i+1}$ 上\Term{action}{作用}平凡。此群又与
\[\prod_{\operatorname{ord}(\alpha)=i+1}U_\alpha\]
等同，因而具有向量结构。
若对 $G$ 的另一个\Term{splitting}{分裂} $(G,T',M',R')$，$B$ 形如 $B_{R'_+}$，证明用这个新\Term{splitting}{分裂}按上述方式构造的群 $U'_i$ 与 $U_i$ 相同，且逐次商上的向量结构也相同。由 5.6.13，存在 $b\in B^u(S)$ 使 $T'=\operatorname{int}(b)T$；待证断言在 $S$ 上是\Term{local-scheme}{局部}的，故可设 $\operatorname{int}(b)$ 诱导的\Term{isomorphism}{同构} $T\xrightarrow\sim T'$，由一个\Term{root-datum}{根数据}\Term{isomorphism}{同构}
\[h:(M',M'^*,R',R'^*)\xrightarrow\sim(M,M^*,R,R^*)\]
\Term{duality}{对偶}而来。显然，$R'_+$ 的根是 $\alpha\circ\operatorname{int}(b)=h(\alpha)$（$\alpha\in R_+$），其\Term{simple-root}{单根}是 $\alpha\circ\operatorname{int}(b)=h(\alpha)$（$\alpha\in\Delta$），故对 $\alpha\in R_+$，有 $\operatorname{ord}(h(\alpha))=\operatorname{ord}(\alpha)$。另一方面，由结构搬移，\Term{vector-group}{向量群} $U'_{h(\alpha)}$ 恰为 $\operatorname{int}(b)U_\alpha$。故 $\operatorname{int}(b)U_i=U'_i$；因 $U_i$ \Term{normal}{正规}，得 $U_i=U'_i$。
同样，由已证结果，\Term{vector-group}{向量群}\Term{isomorphism}{同构}
\[\operatorname{int}(b):U_i/U_{i+1}\xrightarrow\sim U'_i/U'_{i+1}\]
为恒等。
\SourcePage{161}
现在处理一般情形。存在\Term{etale}{平展}\Term{topology}{拓扑}的一个\Term{covering-family}{覆盖族} $\{S_i\to S\}$，并对每个 $i$ 存在一个\Term{splitting}{分裂} $(G_i,T_i,M_i,R_i)$ 及 $R_i$ 的一个\Term{positive-root-system}{正根系} $R_i^+$，使 $B\times_SS_i=B_{R_i^+}$（5.5.5 (iii)）。故对每个 $i$，有一族
\[B_{S_i}=U_{i,0}\supset U_{i,1}\supset\cdots\supset U_{i,j}\supset\cdots\]
及 $U_{i,j}/U_{i,j+1}$ 上的向量结构。由\Term{descent}{下降}，只须证明：对每对 $(i,i')$ 和每个 $j$，有
\[U_{i,j}\times_{S_i}S_{ii'}=U_{i',j}\times_{S_{i'}}S_{ii'},\qquad S_{ii'}=S_i\times_SS_{i'},\]
且商
\[(U_{i,j}/U_{i,j+1})\times_{S_i}S_{ii'},\qquad
(U_{i',j}/U_{i',j+1})\times_{S_{i'}}S_{ii'}\]
上的向量结构相同。若 $S_{ii'}=\varnothing$，显然成立；若 $S_{ii'}\ne\varnothing$，则处于前述情形：$B\times_SS_{ii'}$ 在\Term{splitting}{分裂} $(G_{S_{ii'}},T_i\times_{S_i}S_{ii'},M_i,R_i)$（分别为 $(G_{S_{ii'}},T_{i'}\times_{S_{i'}}S_{ii'},M_{i'},R_{i'})$）中由\Term{positive-root-system}{正根系} $R_i^+$（分别为 $R_{i'}^+$）定义。
\end{proof}
\begin{SGAresult}{推论}{5.9.6}
若 $S$ \Term{affine}{仿射}，则 $H^1(S,B^u)=e$，即 $B^u$ 下的每个\Term{principal-bundle}{主纤维丛}均具有\Term{section}{截面}。
\end{SGAresult}
\begin{proof}
$S$ \Term{resolution}{分解}成一些子\Term{scheme}{概形}的\Term{direct-sum}{直和}，在每个子\Term{scheme}{概形}上 $B^u$ 的\Term{relative-dimension}{相对维数}为常数。故由 5.9.5 (iii)，可设存在 $n$ 使 $U_n=e$。由 TDTE I，B.1.1，\footnote{编者注：指 TDTE I 第 18 页的推论。}有
\[H^1(S,U_i/U_{i+1})=H^1(S,\mathbf W(\mathcal E_i))=0,\]
故 $H^1(S,B^u)=0$。
\end{proof}
\begin{SGAresult}{推论}{5.9.7}
若 $S$ \Term{affine}{仿射}，则 $B$ 具有\Term{maximal-torus}{极大环面}。若 $T$ 为 $B$ 的一个\Term{maximal-torus}{极大环面}，则 $H^1(S,T)=H^1(S,B)$。
\end{SGAresult}
第一个断言立即由 5.9.6 和 5.6.13 得到；第二个以标准方式由此推出。\footnote{编者注：事实上，有\Term{exact-sequence}{正合列} $H^1(S,B^u)\to H^1(S,B)\xrightarrow\pi H^1(S,B/B^u)=H^1(S,T)$，见 [Se64]，I \S 5.5，命题 38 或 [Gi71]，III 命题 3.3.1。而 $H^1(S,T)\to H^1(S,B)$ 是 $\pi$ 的一个\Term{section}{截面}，故 $\pi$ \Term{surjectivity}{满射}；另一方面，由 5.9.6，$H^1(S,B^u)=0$。}
\begin{SGAresult}{推论}{5.9.8}
若 $G$ 为 $S$ 上的\Term{reductive-group}{约化群}，则\Term{canonical}{典范}\Term{morphism}{态射}（参见 5.8.3）$\operatorname{Kil}(G)\to\operatorname{Bor}(G)$ 在每个\Term{affine}{仿射}\Term{open-subset}{开集}上均具有\Term{section}{截面}。
\end{SGAresult}
\begin{SGAresult}{推论}{5.9.9}
在 5.9.5 的条件下，设 $S$ \Term{affine}{仿射}，则存在一个\Term{locally-free}{局部自由} $\mathcal O_S$ 模 $\mathcal E$，使 $B^u$ 作为\Term{scheme}{概形}与 $\mathbf W(\mathcal E)$ 在 $S$ 上\Term{isomorphism}{同构}。
\end{SGAresult}
\SourcePage{162}
\begin{proof}
对 $i$ 归纳证明 $B^u/U_i$ 在 $S$ 上与 $\mathbf W(\mathcal E_0\oplus\cdots\oplus\mathcal E_{i-1})$ \Term{isomorphism}{同构}。$i=0$ 时显然；设 $i\ge1$。$B^u/U_i$ 是以 $X=B^u/U_{i-1}$ 为底、在群 $(U_{i-1}/U_i)_X$ 下的一个\Term{principal-homogeneous-space}{主齐性纤维丛}。由归纳假设，$B^u/U_{i-1}$ \Term{affine}{仿射}，又 $U_{i-1}/U_i=\mathbf W(\mathcal E_{i-1})$，故此\Term{fibre-bundle}{纤维丛}平凡。因此（至少）有一个 $S$ 上\Term{scheme}{概形}\Term{isomorphism}{同构}
\[B^u/U_i\xrightarrow\sim(B^u/U_{i-1})\times_S\mathbf W(\mathcal E_{i-1})
=\mathbf W(\mathcal E_0\oplus\cdots\oplus\mathcal E_{i-1}).\]
由 5.9.5 的条件 (iii) 立即得出结论。
\end{proof}
\begin{SGAresult}{推论}{5.9.10}
设 $S$ 为半\Term{local-scheme}{局部}\Term{scheme}{概形}，$\{s_i\}$ 为其\Term{closed-point}{闭点}，$B$ 为 $S$ 上的\Term{reductive-group}{约化群} $G$ 的一个\Term{borel-subgroup}{Borel 子群}。\Term{canonical}{典范}映射
\[B^u(S)\longrightarrow\prod_i B^u(\operatorname{Spec}\kappa(s_i))\]
为\Term{surjectivity}{满射}。
\end{SGAresult}
\begin{proof}
若 $S=\operatorname{Spec}(A)$，$\kappa(s_i)=A/\mathfrak p_i$，且 $\mathcal E$ 由 $A$ 模 $E$ 给出，则
\[B^u(S)=E\otimes A,\qquad B^u(\operatorname{Spec}\kappa(s_i))=E\otimes_A(A/\mathfrak p_i).\]
断言立即由 $A\to\prod_iA/\mathfrak p_i$ \Term{surjectivity}{满射}得到。
\end{proof}
\SGAsection{5.10}{具有约化纤维的 (R) 型子群}
\begin{SGAresult}{命题}{5.10.1}
设 $(G,T,M,R)$ 为一个 $S$ 上的\Term{split-group}{分裂群}，$R'$ 为 $R$ 的一个 (R) 型\Term{subset}{子集}（5.4.2），$H_{R'}$ 为 $G$ 的相应\Term{subgroup}{子群}。下列条件\Term{equivalence}{等价}：

(i) $H_{R'}$ \Term{reductive-group}{约化}（即\Term{geometric-fibre}{几何纤维}\Term{reduced}{约化}）。

(ii) $R'=-R'$，即 $R'$ 对称。

此外，在这些条件下，$(H_{R'},T,M,R')$ 是 $H_{R'}$ 的一个\Term{splitting}{分裂}；对 $R$ 的每个\Term{positive-root-system}{正根系} $R_+$，$R'_+=R'\cap R_+$ 是 $R'$ 的一个\Term{positive-root-system}{正根系}，且 $B_{R_+}\cap H_{R'}=H_{R'_+}$ 是 $H_{R'}$ 的一个\Term{borel-subgroup}{Borel 子群}，其\Term{unipotent-part}{幺幂部分}为 $U_{R_+}\cap H_{R'}=U_{R'_+}$。
\end{SGAresult}
\begin{proof}
显然 (i)$\Rightarrow$(ii)（只须逐\Term{fiber}{纤维}验证，且 $R'$ 是 $H_{R'}$ 关于 $T$ 的一个\Term{root-system}{根系}）。为证 (ii)$\Rightarrow$(i)，由 5.4.3 注意到对每个 $\alpha\in R'$，有
\[H_{R'}\cap Z_\alpha=\operatorname{Centr}_{H_{R'}}(T_\alpha)=Z_\alpha,\]
再应用第 XIX 讲 1.12 的判据。
若 $R_+$ 为 $R$ 的一个\Term{positive-root-system}{正根系}，则 $R'_+=R_+\cap R'$ 显然是 $R'$ 的一个\Term{closed-root-subset}{闭子集}，满足 $R'_+\cup-R'_+=R'$ 及 $R'_+\cap-R'_+=\varnothing$，故为 $R'$ 的一个\Term{positive-root-system}{正根系}。其余两个断言分别由 5.6.1 (vi) 和 5.6.7 (i) 得到。
\end{proof}
\SourcePage{163}
\begin{SGAresult}{推论}{5.10.2}
设 $S$ 为\Term{scheme}{概形}，$G$ 为 $S$ 上的\Term{reductive-group}{约化群}，$H$ 为一个\Term{reductive-group}{约化群}子\Term{scheme}{概形}，使对每个 $s\in S$，$G_s,H_s$ 的\Term{reductive-rank}{约化秩}相同。则 $H$ 在 $G$ 中闭，$\operatorname{Norm}_G(H)$ 在 $S$ 上\Term{smooth}{光滑}，$\operatorname{Norm}_G(H)/H$ 可由一个 $S$ 上的\Term{finite}{有限}\Term{etale}{平展}\Term{scheme}{概形}表示。
若 $T$ 为 $H$ 的一个\Term{maximal-torus}{极大环面}，$B$ 为 $G$ 中包含 $T$ 的一个\Term{borel-subgroup}{Borel 子群}，则 $B\cap H$ 为 $H$ 的一个\Term{borel-subgroup}{Borel 子群}，其\Term{unipotent-part}{幺幂部分}为 $(B\cap H)^u=B^u\cap H$。
\end{SGAresult}
\begin{proof}
前几个断言立即由 5.3.10 和 5.3.18 以及 $G,H$ 的\Term{weyl-group}{Weyl 群}\Term{finite}{有限}（第 XIX 讲 2.5）得到。其余断言对\Term{etale}{平展}\Term{topology}{拓扑}是\Term{local-scheme}{局部}的，归结为 5.10.1 中考察的情形。
\end{proof}
\begin{SGAresult}{命题}{5.10.3}
设 $S$ 为\Term{scheme}{概形}，$G$ 为 $S$ 上的\Term{reductive-group}{约化群}。

a) 若 $Q$ 为 $G$ 的一个\Term{torus}{环面}，则 $\operatorname{Centr}_G(Q)$ 为 $G$ 的一个具有\Term{reduced}{约化}\Term{fiber}{纤维}的 (R) 型\Term{subgroup}{子群}。若 $Q\subset Q'$ 是 $G$ 的两个\Term{torus}{环面}，则 $\operatorname{Centr}_G(Q)\supset\operatorname{Centr}_G(Q')$。

b) 若 $H$ 为 $G$ 的一个具有\Term{reduced}{约化}\Term{fiber}{纤维}的 (R) 型\Term{subgroup}{子群}，则 $\operatorname{rad}(H)$（4.3.6）是 $G$ 的一个\Term{torus}{环面}。若 $H\subset H'$ 为 $G$ 的两个具有\Term{reduced}{约化}\Term{fiber}{纤维}的 (R) 型\Term{subgroup}{子群}，则 $\operatorname{rad}(H)\supset\operatorname{rad}(H')$。

c) 若 $Q$ 为 $G$ 的一个\Term{torus}{环面}，则
\[\operatorname{rad}(\operatorname{Centr}_G(Q))\supset Q,\qquad
\operatorname{Centr}_G(\operatorname{rad}(\operatorname{Centr}_G(Q)))=\operatorname{Centr}_G(Q).\]
d) 若 $H$ 为 $G$ 的一个具有\Term{reduced}{约化}\Term{fiber}{纤维}的 (R) 型\Term{subgroup}{子群}，则
\[\operatorname{Centr}_G(\operatorname{rad}(H))\supset H,\qquad
\operatorname{rad}(\operatorname{Centr}_G(\operatorname{rad}(H)))=\operatorname{rad}(H).\]
\end{SGAresult}
\begin{proof}
a) 立即由第 XIX 讲 2.8 得到。为证 b)，只须注意 $\operatorname{rad}(H')\subset H$，因为 $H$ 对 (fpqc) \Term{local-scheme}{局部}包含 $G$ 的一个\Term{maximal-torus}{极大环面}，因而也包含 $H'$ 的一个\Term{maximal-torus}{极大环面}。c)（分别为 d)）的第一个断言显然，第二个由通常的论证得到。
\end{proof}
此命题引出下述定义。
\begin{SGAresult}{定义}{5.10.4}
设 $S$ 为\Term{scheme}{概形}，$G$ 为 $S$ 上的\Term{reductive-group}{约化群}，$H$ 为 $G$ 的一个 (R) 型\Term{reductive-group}{约化子群}，$Q$ 为 $G$ 的一个\Term{subtorus}{子环面}。\footnote{编者注：修改了原版，以“\Term{c-critical-torus}{C-临界环面}”代替“\Term{critical-torus}{临界环面}”，避免后续引文中的混淆（参见第 XXVI 讲 3.9）。还详细补充了 5.10.5 的陈述，并加入注记 5.10.5.1。}

1) 若 $H$ 为其\Term{radical}{根基}的\Term{centralizer}{中心化子}，则称 $H$ 为\Term{critical-subgroup}{临界子群}（sous-groupe critique）。

2) 若 $Q$ 为其\Term{centralizer}{中心化子}的\Term{radical}{根基}，则称 $Q$ 为\Term{c-critical-torus}{C-临界环面}（tore C-critique）。
\end{SGAresult}
由命题 5.10.3 得到：
\begin{SGAresult}{推论}{5.10.5}
(i) 对 $G$ 的每个\Term{subtorus}{子环面} $Q$，$\operatorname{Centr}_G(Q)$ 是\Term{critical-subgroup}{临界子群}。

(ii) 对 $G$ 的每个具有\Term{reduced}{约化}\Term{fiber}{纤维}的 (R) 型\Term{subgroup}{子群} $H$，$\operatorname{rad}(H)$ 是 $G$ 的一个\Term{c-critical-torus}{C-临界环面}。

(iii) 映射 $Q\mapsto\operatorname{Centr}_G(Q)$、$H\mapsto\operatorname{rad}(H)$，
\SourcePage{164}
在 $G$ 的\Term{c-critical-torus}{C-临界环面}的\Term{set}{集合}与其 (R) 型\Term{critical-subgroup}{临界约化子群}的\Term{set}{集合}之间给出互逆的\Term{bijection}{双射}。

(iv) 若 $Q$ 为 $G$ 的一个\Term{torus}{环面}，则 $\operatorname{rad}(\operatorname{Centr}_G(Q))$ 是 $G$ 中包含 $Q$ 的最小\Term{c-critical-torus}{C-临界环面}。

(v) 若 $H$ 为 $G$ 的一个 (R) 型\Term{reductive-group}{约化子群}，则 $\operatorname{Centr}_G(\operatorname{rad}(H))$ 是 $G$ 中包含 $H$ 的最小 (R) 型\Term{critical-subgroup}{临界约化子群}。
\end{SGAresult}
\textbf{注记 5.10.5.1。}（编者增补，参见脚注 54。）1) $G$ 的一个\Term{torus}{环面} $T$ 是 $G$ 的\Term{critical-subgroup}{临界子群}，当且仅当它是一个\Term{maximal-torus}{极大环面}。2) 下文中，“\Term{c-critical-torus}{临界环面}”指“\Term{c-critical-torus}{C-临界环面}”。
\begin{SGAresult}{命题}{5.10.6}
设 $(G,T,M,R)$ 为一个 $S$ 上的\Term{split-group}{分裂群}，$R'$ 为 $R$ 的一个\Term{subset}{子集}。下列条件\Term{equivalence}{等价}：

(i) $R'$ 为 (R) 型，且 $H_{R'}$ \Term{reductive-group}{约化}且\Term{critical-subgroup}{临界}。

(ii) 存在 $R$ 的一个\Term{simple-root}{单根}系 $\Delta$ 及 $\Delta$ 的一个\Term{subset}{子集} $\Delta'$，使 $R'$ 为 $R$ 中所有\Term{representable}{可表}为 $\Delta'$ 中元素的线性组合的元素的\Term{set}{集合}。

(iii) $R'$ 为\Term{closed-root-subset}{闭}、对称\Term{subset}{子集}，且 $R'$ 的任意\Term{simple-root}{单根}系均为 $R$ 的某个\Term{simple-root}{单根}系与 $R'$ 的交。
\end{SGAresult}
\begin{proof}
由第 XXI 讲 3.4.8，(ii)、(iii) \Term{equivalence}{等价}，且都\Term{equivalence}{等价}于 $R'$ 是 $R$ 与 $M\otimes\mathbb Q$ 的某个 $\mathbb Q$ 向量子\Term{space}{空间}的交。最后这个条件由 (i) 推出：若 $H_{R'}=\operatorname{Centr}_G(Q)$，则 $R'$ 是 $R$ 中在 $Q$ 上为零的元素的\Term{set}{集合}（第 II 讲 5.2.3 (ii)）。最后，此条件推出 (i)：因为 $\operatorname{rad}(H_{R'})$ 是 $\bigcap_{\alpha\in R'}\operatorname{Ker}(\alpha)$ 的\Term{maximal-torus}{极大环面}，故 $\operatorname{Centr}_G(\operatorname{rad}(H_{R'}))$ 恰为 $H_{R''}$，其中 $R''$ 是 $R$ 与 $R'$ 生成的向量子\Term{space}{空间}的交。
\end{proof}
\textbf{5.10.7。} 汇总前述若干结果。设 $(G,T,M,R)$ 为一个 $S$ 上的\Term{split-group}{分裂群}，$\Delta$ 为 $R$ 的一个\Term{simple-root}{单根}系，$R_+$ 为相应\Term{positive-root-system}{正根系}；取 $\Delta$ 的一个\Term{subset}{子集} $\Delta'$，用 $R'$ 表示 $R$ 中所有\Term{representable}{可表}为 $\Delta'$ 中元素的线性组合的元素的\Term{set}{集合}，并置 $R'_+=R'\cap R_+$。令 $T_{\Delta'}$ 为 $\bigcap_{\alpha\in\Delta'}\operatorname{Ker}(\alpha)$ 的\Term{maximal-torus}{极大环面}，$Z_{\Delta'}=\operatorname{Centr}_G(T_{\Delta'})$。
则 $Z_{\Delta'}$ 为 $G$ 的一个\Term{reductive-group}{约化子群}，\Term{radical}{根基}为 $T_{\Delta'}$；$(Z_{\Delta'},T,M,R')$ 是一个 $S$ 上的\Term{split-group}{分裂群}；$B_{R_+}\cap Z_{\Delta'}$ 是由\Term{positive-root-system}{正根系} $R'_+$（或\Term{simple-root}{单根}系 $\Delta'$）定义的 $Z_{\Delta'}$ 的\Term{borel-subgroup}{Borel 子群}，其\Term{unipotent-part}{幺幂部分}为 $U_{R_+}\cap Z_{\Delta'}=U_{R'_+}$。

\textbf{注记 5.10.8。} 在 5.10.4 的条件下，设 $Q$ 为 $G$ 的一个\Term{c-critical-torus}{临界环面}，$L=\operatorname{Centr}_G(Q)$ 为其\Term{centralizer}{中心化子}。因 $Q=\operatorname{rad}(L)$，$Q$ 是 $L$ 的\Term{characteristic-subgroup}{特征子群}；立即推出 $\operatorname{Norm}_G(L)=\operatorname{Norm}_G(Q)$，从而
\[\operatorname{Norm}_G(L)/L=\operatorname{Norm}_G(Q)/\operatorname{Centr}_G(Q)=W_G(Q).\]
由 5.10.2，得到：
\SourcePage{165}
\begin{SGAresult}{命题}{5.10.9}
设 $S$ 为\Term{scheme}{概形}，$G$ 为 $S$ 上的\Term{reductive-group}{约化群}，$Q$ 为 $G$ 的一个\Term{c-critical-torus}{临界环面}。\Term{weyl-group}{Weyl 群} $W_G(Q)$ 在 $S$ 上（\Term{etale}{平展}且）\Term{finite}{有限}。
\end{SGAresult}
\textbf{注记 5.10.10。} 在 5.10.7 的条件下，可将 $W_G(T_{\Delta'})=\operatorname{Norm}_G(Z_{\Delta'})/Z_{\Delta'}$ 明确写出。它是\Term{quotient-group}{商群} $W_1/W_2$ 对应的\Term{constant-group}{常值群}，其中 $W_1$ 为 $W$ 中\Term{normalization}{正规化} $\Delta'$ 生成的 $M$ 的\Term{subgroup}{子群}的那些元素组成的\Term{subgroup}{子群}，$W_2$ 为 $W$ 中由 $s_\alpha$（$\alpha\in\Delta'$）生成的\Term{subgroup}{子群}。
\SGAsection{5.11}{(RC) 型子群}
\begin{SGAresult}{定义}{5.11.1}
设 $S$ 为\Term{scheme}{概形}，$G$ 为 $S$ 上的\Term{reductive-group}{约化群}。若 $G$ 的一个群子\Term{scheme}{概形} $H$ 为 (R) 型，即满足以下两个条件（5.2.1）：

(i) $H$ 在 $S$ 上\Term{smooth}{光滑}，\Term{fiber}{纤维}\Term{connected}{连通}；

(ii) 对每个 $s\in S$，$H_s$ 包含 $G_s$ 的一个\Term{maximal-torus}{极大环面}；

并且还满足下述条件，则称其为 (RC) 型：

(iii) 对每个 $s\in S$ 和 $H_s$ 的每个\Term{maximal-torus}{极大环面} $T$，$H_s$ 关于 $T$ 的\Term{root}{根}的\Term{set}{集合}为 $G_s$ 关于 $T$ 的全部\Term{root}{根}的\Term{set}{集合}的一个\Term{closed-root-subset}{闭子集}。
\end{SGAresult}
\textbf{注记 5.11.2。} 如 5.4.8 已指出，当 $6$ 在 $S$ 上\Term{invertible}{可逆}时，条件 (iii) 由其他条件推出。\footnote{编者注：即当 $S$ 的每个\Term{residue-characteristic}{剩余特征}都大于 $3$ 时。}
\begin{SGAresult}{引理}{5.11.3}
设 $(G,T,M,R)$ 为一个 $S$ 上的\Term{split-group}{分裂群}，$R'$ 为 $R$ 的一个\Term{closed-root-subset}{闭子集}。置
\[R_1=\{\alpha\in R':-\alpha\in R'\},\qquad R_2=\{\alpha\in R':-\alpha\notin R'\}.\]
则 $R_1,R_2$ 均为\Term{closed-root-subset}{闭子集}。考虑\Term{smooth}{光滑}且\Term{fiber}{纤维}\Term{connected}{连通}的群 $H_{R'}$、$H_{R_1}$、$U_{R_2}$（5.4.7 和 5.6.5）。

(i) $U_{R_2}$ 在 $H_{R'}$ 中\Term{normal}{正规}，且 $H_{R'}$ 是 $U_{R_2}$ 与 $H_{R_1}$ 的\Term{semidirect-product}{半直积}。

(ii) $H_{R_1}$ \Term{reductive-group}{约化}，$U_{R_2}$ 的\Term{geometric-fibre}{几何纤维}\Term{connected}{连通}且\Term{unipotent}{幺幂}；$H_{R'}$ 中每个在 $S$ 上\Term{smooth}{光滑}、\Term{geometric-fibre}{几何纤维}\Term{connected}{连通}且\Term{unipotent}{幺幂}的\Term{normal-subgroup}{正规子群}都包含于 $U_{R_2}$；而 $H_{R'}$ 中每个包含 $T$ 的\Term{reductive-group}{约化子群}都包含于 $H_{R_1}$。

(iii) $U_{R_2}\cap\operatorname{Norm}_G(H_{R_1})=e$。
\end{SGAresult}
\begin{proof}
先由 5.6.7 (i) 得到 (iii)。(i) 的第一个断言由 5.6.7 (ii) 得到。因由 (iii) 有 $U_{R_2}\cap H_{R_1}=e$，\Term{semidirect-product}{半直积} $H_{R_1}\cdot U_{R_2}$ 为 $H_{R'}$ 的一个\Term{subgroup}{子群}；二者均为 $G$ 中包含 $T$ 的 (R) 型\Term{subgroup}{子群}，且有相同的\Term{lie-algebra}{Lie 代数} $\mathfrak g_{R'}$；由 5.3.5，它们相同，(i) 得证。
再证明 (ii)：前两个断言就是 5.10.1 和 5.6.5。设 $U$ 为 $H_{R'}$ 的一个\Term{smooth}{光滑}、\Term{finite-presentation}{有限呈示}、
\SourcePage{166}
\Term{normal}{正规}（因而由 $T$ \Term{normalization}{正规化}）且\Term{geometric-fibre}{几何纤维}\Term{connected}{连通}\Term{unipotent}{幺幂}的群子\Term{scheme}{概形}。由 5.6.12，在 $S$ 上\Term{local-scheme}{局部}有 $U=U_{R''}$，其中 $R''$ 为 $R'$ 的一个\Term{subset}{子集}，满足 $R''\cap-R''=\varnothing$。若 $U\not\subset U_{R_2}$，则 $R''\not\subset R_2$，故存在 $\alpha\in R''$ 使 $-\alpha\in R'$。于是 $Z_\alpha\subset H_{R'}$（5.4.3），故 $Z_\alpha$ \Term{normalization}{正规化} $U$。但 $U$ 包含 $U_\alpha$，且 $Z_\alpha$ 有一个\Term{section}{截面} $w$ 使 $\operatorname{int}(w)U_\alpha=U_{-\alpha}$；这推出 $-\alpha\in R''$，与 $R''\cap-R''=\varnothing$ 矛盾。
最后，若 $L$ 为 $H_{R'}$ 中包含 $T$ 的一个\Term{reductive-group}{约化子群}，则在 $S$ 上\Term{local-scheme}{局部}有 $L=H_{R'''}$，其中 $R'''$ 为包含于 $R'$ 的对称\Term{subset}{子集}，故包含于 $R_1$。
\end{proof}
\begin{SGAresult}{命题}{5.11.4}
设 $S$ 为\Term{scheme}{概形}，$G$ 为 $S$ 上的\Term{reductive-group}{约化群}，$H$ 为 $G$ 的一个 (RC) 型群子\Term{scheme}{概形}。

(i) $H$ 在 $G$ 中闭，$\operatorname{Norm}_G(H)/H$ 可由一个 $S$ 上的\Term{finite}{有限}\Term{etale}{平展}群\Term{scheme}{概形}表示。

(ii) $H$ 具有一个最大的在 $S$ 上\Term{smooth}{光滑}且\Term{finite-presentation}{有限呈示}、\Term{geometric-fibre}{几何纤维}\Term{connected}{连通}且\Term{unipotent}{幺幂}的\Term{normal-subgroup}{正规子群}\Term{scheme}{概形}；称之为 $H$ 的\Term{unipotent-radical}{幺幂根基}，记为 $\operatorname{rad}_u(H)$。商\Term{sheaf}{层} $H/\operatorname{rad}_u(H)$ 可由一个 $S$ 上的\Term{reductive-group}{约化群}表示。

(iii) 若 $T$ 为 $H$ 的一个\Term{maximal-torus}{极大环面}，则 $H$ 具有一个包含 $T$ 的 (RC) 型\Term{reductive-group}{约化子群} $L$，具有下列两个性质：

(a) $H$ 中每个包含 $T$ 的\Term{reductive-group}{约化子群}均包含于 $L$。

(b) $H$ 是\Term{semidirect-product}{半直积} $H=L\cdot\operatorname{rad}_u(H)$，即\Term{canonical}{典范}\Term{morphism}{态射} $L\to H/\operatorname{rad}_u(H)$ 为\Term{isomorphism}{同构}。

此外，$L$ 是 $H$ 中包含 $T$、并满足上述两个条件任意一个的唯一\Term{reductive-group}{约化子群}。最后有
\[\operatorname{Norm}_H(L)=L,\qquad\operatorname{Norm}_H(T)=\operatorname{Norm}_L(T),\qquad W_H(T)=W_L(T);\]
特别地，$W_H(T)$ 在 $S$ 上\Term{finite}{有限}。
\end{SGAresult}
\begin{proof}
先注意 (i) 对\Term{etale}{平展}\Term{topology}{拓扑}是\Term{local-scheme}{局部}的。因此由推论 5.3.18，(i) 是 (iii) 最后一个断言的推论。
(ii) 的断言对\Term{etale}{平展}\Term{topology}{拓扑}是\Term{local-scheme}{局部}的，故可设处于 5.11.3 的情形，立即由其中的 (i)、(ii) 得到结论。
由其所含的唯一性断言，(iii) 对\Term{etale}{平展}\Term{topology}{拓扑}也是\Term{local-scheme}{局部}的，仍可归结为 5.11.3 的情形，其中已验证性质 (a)、(b)。满足 (a) 的 $L$ 的唯一性显然；由 (a)，满足 (b) 的 $L$ 的唯一性也显然。等式 $\operatorname{Norm}_H(L)=L$ 就是 5.11.3 (iii)；若 $H$ 的一个\Term{section}{截面}\Term{normalization}{正规化} $T$，则由 $L$ 的唯一性，它也\Term{normalization}{正规化} $L$，故由刚才的结果，它是 $L$ 的一个\Term{section}{截面}，这证明了第二个等式；第三个随即显然。
\end{proof}
\begin{SGAresult}{命题}{5.11.5}
设 $S$ 为\Term{scheme}{概形}，$G$ 为 $S$ 上的\Term{reductive-group}{约化群}，$\mathcal H_c$ 为 $G$ 的 (RC) 型\Term{subgroup}{子群}\Term{functor}{函子}，它是 5.8.1 的\Term{functor}{函子} $\mathcal H$ 的一个子\Term{functor}{函子}。

(i) $\mathcal H_c$ 可由 $\mathcal H$ 的一个在 $S$ 上\Term{smooth}{光滑}、\Term{quasi-projective}{拟射影}且\Term{finite-presentation}{有限呈示}的开子\Term{scheme}{概形}表示。
\SourcePage{167}
(ii) 存在一个 $S$ 上的\Term{finite}{有限}\Term{etale}{平展}\Term{scheme}{概形} $\mathcal{Cl}_c$ 及一个\Term{smooth}{光滑}、\Term{quasi-projective}{拟射影}、\Term{finite-presentation}{有限呈示}、\Term{surjectivity}{满射}且\Term{geometric-fibre}{几何纤维}\Term{connected}{连通}的\Term{morphism}{态射}
\[\mathrm{cl}:\mathcal H_c\longrightarrow\mathcal{Cl}_c,\]
满足：对每个 $S'\to S$ 及 $H,H'\in\mathcal H_c(S')$，$\mathrm{cl}(H)=\mathrm{cl}(H')$ 当且仅当 $H,H'$ 在 $G$ 中对\Term{etale}{平展}\Term{topology}{拓扑}\Term{local-scheme}{局部}\Term{conjugacy}{共轭}（由 5.3.11，这又\Term{equivalence}{等价}于对每个 $s\in S$，$H_s,H'_s$ 由 $G(\overline s)$ 的某个元素\Term{conjugacy}{共轭}）。

(iii) 上述条件唯一确定 $\mathcal{Cl}_c$ 和 $\mathrm{cl}$（精确到唯一\Term{isomorphism}{同构}）。

(iv) 若 $(G,T,M,R)$ 为 $G$ 的一个\Term{splitting}{分裂}，令 $E$ 为 $R$ 的\Term{closed-root-subset}{闭子集}模 $W$ \Term{conjugacy}{共轭}的类的\Term{set}{集合}；则存在\Term{isomorphism}{同构} $\mathcal{Cl}_c\xrightarrow\sim E_S$，使对 $R$ 的任意\Term{closed-root-subset}{闭子集} $R'$，$\mathrm{cl}(H_{R'})$ 对应于 $R'$ 在 $E_S(S)=\operatorname{Hom}_{\mathrm{loc.const.}}(S,E)$ 中的\Term{canonical}{典范}像。
\end{SGAresult}
\begin{proof}
首先，$\mathcal H_c$ 显然是\Term{etale}{平展}\Term{topology}{拓扑}的\Term{sheaf}{层}，且 (ii) 推出 $\mathcal{Cl}_c$ 就是 $\mathcal H_c$ 模由\Term{conjugacy}{共轭}定义的\Term{equivalence-relation}{等价关系}的商\Term{sheaf}{层}。这首先推出 (iii)，也说明只须对\Term{etale}{平展}\Term{topology}{拓扑}\Term{local-scheme}{局部}验证 (i)、(ii)。故归结到 (iv) 的情形；先构造一个\Term{morphism}{态射} $f:\mathcal H_c\to E_S$。
只须构造一个关于 $S$ \Term{functor}{函子}性的映射 $\mathcal H_c(S)\to E_S(S)$。设 $H$ 为 $G$ 的一个 (RC) 型\Term{subgroup}{子群}；因 $H$ 对\Term{etale}{平展}\Term{topology}{拓扑}\Term{local-scheme}{局部}具有\Term{maximal-torus}{极大环面}，而 $G$ 的\Term{maximal-torus}{极大环面}对\Term{etale}{平展}\Term{topology}{拓扑}\Term{local-scheme}{局部}\Term{conjugacy}{共轭}，故存在一个\Term{covering-family}{覆盖族} $\{S_i\to S\}$，且对每个 $i$ 有 $g_i\in G(S_i)$ 及 $R$ 的一个\Term{closed-root-subset}{闭子集} $R_i$，使
\[\operatorname{int}(g_i)(H\times_SS_i)=H_{R_i}\times_SS_i.\]
每个 $R_i$ 定义 $E_{S_i}$ 的一个\Term{section}{截面} $\eta_i$，即 $E_S(S_i)$ 的一个元素；只须证明族 $(\eta_i)$ 来自 $E_S$ 在 $S$ 上的一个\Term{section}{截面} $\eta=f(H)$，且它仅依赖于 $H$。
为此，归结到证明 $H_{R'}$ 与 $H_{R''}$ 对\Term{etale}{平展}\Term{topology}{拓扑}\Term{local-scheme}{局部}\Term{conjugacy}{共轭}，当且仅当 $R',R''$ 由 \Term{weyl-group}{Weyl 群} $W$ 的一个元素\Term{conjugacy}{共轭}；这显然成立。
对每个 $\eta\in E$，存在 $H_0\in\mathcal H_c(S)$ 使 $f(H_0)=\eta$：只须取 $H_0=H_{R'}$，其中 $R'$ 为在 $E$ 中像为 $\eta$ 的 $R$ 的一个\Term{closed-root-subset}{闭子集}。若 $H\in\mathcal H_0(S')$、$S'\to S$，则 $H$ 对\Term{etale}{平展}\Term{topology}{拓扑}\Term{local-scheme}{局部}与 $H_0$ \Term{conjugacy}{共轭}，当且仅当 $f(H)=\eta$（由前面的论证立即看出）。故 $f^{-1}(\eta)$ 与商 $G/\operatorname{Norm}_G(H_0)$ 等同，而由 5.8.2，此商为 $\mathcal H$ 的一个在 $S$ 上\Term{smooth}{光滑}、\Term{quasi-projective}{拟射影}、\Term{finite-presentation}{有限呈示}、\Term{fiber}{纤维}\Term{connected-nonempty}{连通非空}的\Term{open-subset}{开集}。因 $E_S$ 是对应于各 $\eta\in E$ 的\Term{section}{截面}的像开子\Term{scheme}{概形}的和，$\mathcal H_c$ 与所有 $f^{-1}(\eta)$（$\eta\in E$）的和等同，(i)、(ii) 得证。最后，由构造可知 (iv) 成立。
\end{proof}
\begin{SGAresult}{推论}{5.11.6}
若 $u\in\mathcal{Cl}_c(S')$、$S'\to S$，则 $\mathrm{cl}^{-1}(u)$ 为一个\Term{smooth}{光滑}、\Term{quasi-projective}{拟射影}、\Term{finite-presentation}{有限呈示}且\Term{fiber}{纤维}\Term{connected-nonempty}{连通非空}的 $S'$ 上\Term{scheme}{概形}；它是 $\mathcal H_c$ 的一个\Term{open-subset}{开集}，
\SourcePage{168}
并且是 $G_{S'}$ 下的一个“齐性”\Term{scheme}{概形}（\Term{action}{作用}由\Term{inner-automorphism}{内自同构}给出）。特别地，若 $H\in\mathrm{cl}^{-1}(u)(S')$，则由 $g\mapsto\operatorname{int}(g)H$ 定义的\Term{morphism}{态射} $G_{S'}\to(\mathcal H_c)_{S'}$ 将 $G_{S'}/\operatorname{Norm}_{G_{S'}}(H)$ 与 $\mathrm{cl}^{-1}(u)$ 等同。
\end{SGAresult}
\textbf{例 5.11.7。} 特别地，$\mathcal{Cl}_c$ 有两个\Term{canonical}{典范}\Term{section}{截面} $u_t,u_b$，分别对应于\Term{maximal-torus}{极大环面}（$R'=\varnothing$）及\Term{borel-subgroup}{Borel 子群}（$R'$ 为\Term{positive-root-system}{正根系}）。$S$ 上的\Term{scheme}{概形} $\mathrm{cl}^{-1}(u_t)$、$\mathrm{cl}^{-1}(u_b)$ 就是 5.8.3 引入的 $S$ 上\Term{scheme}{概形} $\operatorname{Tor}(G)$、$\operatorname{Bor}(G)$。第 XXVI 讲还将给出其他例子。

\textbf{注记 5.11.8。} 可以构造一个在 $S$ 上\Term{finite-presentation}{有限呈示}且\Term{unramified}{非分歧}的 $S$ 上\Term{scheme}{概形} $\mathcal{Cl}$，以及一个\Term{smooth}{光滑}\Term{surjectivity}{满射} $\mathcal H\to\mathcal{Cl}$，其\Term{geometric-fibre}{几何纤维}\Term{connected}{连通}，并具有类似于 5.11.5 (ii)、(iii) 的性质。
\SGAsection{6}{导群}
\SGAsection{6.1}{预备知识}
本小节固定一个\Term{scheme}{概形} $S$、一个 $S$ 上的\Term{split-group}{分裂群} $(G,T,M,R)$、$R$ 的一个\Term{positive-root-system}{正根系} $R_+$，并记
\[B=B_{R_+},\qquad B^-=B_{R_-},\qquad U=B^u,\qquad U^-=(B^-)^u,\]
\[\Omega=\Omega_{R_+}=U^-\cdot T\cdot U.\]
\textbf{6.1.1。} 用 $T'$ 表示 $T$ 中“族 $\alpha^*$（$\alpha\in R$）的像”这一\Term{subtorus}{子环面}；换言之，$T'$ 是由
\[(z_\alpha)_{\alpha\in R}\longmapsto\prod_{\alpha\in R}\alpha^*(z_\alpha)\]
定义的群\Term{morphism}{态射} $\mathbb G_{m,S}^R\to T$ 的像。若 $\Delta$ 表示 $R_+$ 的\Term{simple-root}{单根}\Term{set}{集合}，立即可见，以同样方式定义的\Term{morphism}{态射}
\[\mathbb G_{m,S}^\Delta\longrightarrow T'\]
为\Term{surjectivity}{满射}且核\Term{finite}{有限}。若将 $T$ 与 $D_S(M)$ 等同，则 $T'$ 与 $D_S(M/N)$ 等同，其中
\[N=M\cap\mathscr V(R^*)^\perp\]
（$\mathscr V(R^*)^\perp$ 表示在 $V$ 与 $V^*$ 的\Term{duality}{对偶}下，$\mathscr V(R^*)$ 的\Term{orthogonality}{正交}\Term{space}{空间}）。
\begin{SGAresult}{引理}{6.1.2}
由 $T$ 中乘法定义的\Term{morphism}{态射} $\operatorname{rad}(G)\times_ST'\to T$ 是\Term{isogeny}{同源}（参见 4.2.9）。
\end{SGAresult}
\begin{proof}
\Term{canonical}{典范}\Term{morphism}{态射} $\operatorname{rad}(T)\to T/T'$ 由\Term{abelian-group}{交换群}\Term{morphism}{态射}
\[M\cap\mathscr V(R^*)^\perp\longrightarrow M/(M\cap\mathscr V(R))\]
\Term{duality}{对偶}而来，立即可见后者\Term{injectivity}{单射}且\Term{cokernel}{余核}\Term{finite}{有限}（参见第 XXI 讲 6.3）。
\end{proof}
\SourcePage{169}
\begin{SGAresult}{定义}{6.1.3}
置 $\Omega'=U^-\cdot T'\cdot U$；它是 $\Omega=U^-\cdot T\cdot U$ 的一个闭子\Term{scheme}{概形}。
\end{SGAresult}
\begin{SGAresult}{引理}{6.1.4}
设 $\alpha$ 为一个\Term{simple-root}{单根}，$w_\alpha\in\operatorname{Norm}_G(T)(S)$ \Term{lifting}{提升} $s_\alpha$。则
\[\operatorname{int}(w_\alpha)\Omega'\cap\Omega\subset\Omega'.\]
\end{SGAresult}
\begin{proof}
只须证明：若 $g\in\Omega'(S)$ 且 $\operatorname{int}(w_\alpha)g\in\Omega(S)$，则 $\operatorname{int}(w_\alpha)g\in\Omega'(S)$。由 5.6.8，写成
\[g=a\exp_{-\alpha}(Y)t\exp_\alpha(X)b,\]
其中 $a\in U_{-\widehat\alpha}(S)$、$Y\in\Gamma(S,\mathfrak g^{-\alpha})$、$t\in T'(S)$、$X\in\Gamma(S,\mathfrak g^\alpha)$、$b\in U_{\widehat\alpha}(S)$。于是
\[\operatorname{int}(w_\alpha)g=\operatorname{int}(w_\alpha)a\cdot
\operatorname{int}(w_\alpha)\bigl(\exp_{-\alpha}(Y)t\exp_\alpha(X)\bigr)\cdot\operatorname{int}(w_\alpha)b.\]
由 5.6.8 (iv)，有
\[\operatorname{int}(w_\alpha)a\in U_{-\widehat\alpha}(S),\qquad
\operatorname{int}(w_\alpha)b\in U_{\widehat\alpha}(S).\]
置 $h=\exp_{-\alpha}(Y)t\exp_\alpha(X)$，得到
\[\begin{aligned}
\operatorname{int}(w_\alpha)g\in\Omega(S)&\Longleftrightarrow\operatorname{int}(w_\alpha)h\in\Omega(S),\\
\operatorname{int}(w_\alpha)g\in\Omega'(S)&\Longleftrightarrow\operatorname{int}(w_\alpha)h\in\Omega'(S).
\end{aligned}\]
故归结到 $g=h$。由 4.1.12，
\[Z_\alpha\cap\Omega=U_{-\alpha}\cdot T\cdot U_\alpha,\qquad
Z_\alpha\cap\Omega'=U_{-\alpha}\cdot T'\cdot U_\alpha,\]
所以归结到证明
\[\operatorname{int}(w_\alpha)h\in(U_{-\alpha}\cdot T\cdot U_\alpha)(S)
\Longrightarrow\operatorname{int}(w_\alpha)h\in(U_{-\alpha}\cdot T'\cdot U_\alpha)(S).\]
这立即由第 XX 讲 3.12 得到；那里说明 $\operatorname{int}(w_\alpha)h$ 在 $T$ 上的\Term{set-component}{分量}形如 $t\cdot\alpha^*(z)\in T'(S)$。
\end{proof}
\begin{SGAresult}{引理}{6.1.5}
对每个 $w\in\operatorname{Norm}_G(T)(S)$，存在 $G$ 的一个包含单位\Term{section}{截面}的\Term{open-subset}{开集} $V_w$，使 $\operatorname{int}(w)\Omega'\cap V_w\subset\Omega'$。
\end{SGAresult}
\begin{proof}
对每个\Term{simple-root}{单根} $\alpha$ 取\Term{lifting}{提升} $s_\alpha$ 的 $n_\alpha\in\operatorname{Norm}_G(T)(S)$。对每个 $s\in S$，存在包含 $s$ 的 $S$ 的一个\Term{open-subset}{开集} $V$、$t\in T(V)$ 及 $V$ 上的关系
\[w=n_{\alpha_1}\cdots n_{\alpha_p}t,\qquad\alpha_i\text{ 均为单根}.\]
显然只须在 $V=S$ 时证明；对 $p$ 归纳。若 $p=0$，则 $w\in T(S)$，取 $V_w=G$；故设 $w=n_\alpha w'$，且 $w'$ 满足引理结论，即存在 $G$ 的一个包含单位\Term{section}{截面}的\Term{open-subset}{开集} $V_{w'}$，使 $\operatorname{int}(w')\Omega'\cap V_{w'}\subset\Omega'$。由 6.1.4，
\[\begin{split}
\operatorname{int}(w)\Omega'\cap\operatorname{int}(n_\alpha)V_{w'}\cap\Omega
&=\operatorname{int}(n_\alpha)\bigl(\operatorname{int}(w')\Omega'\cap V_{w'}\bigr)\cap\Omega\\
&\subset\operatorname{int}(n_\alpha)\Omega'\cap\Omega\subset\Omega'.
\end{split}\]
故取 $V_w=\operatorname{int}(n_\alpha)V_{w'}\cap\Omega$，证明完成。
\end{proof}
\SourcePage{170}
\begin{SGAresult}{引理}{6.1.6}
存在 $G$ 的一个包含单位\Term{section}{截面}的\Term{open-subset}{开集} $V_0$，使对每个 $S'\to S$，有
\[U(S')U^-(S')\cap V_0(S')\subset\Omega'(S').\]
\end{SGAresult}
\begin{proof}
取 $n_0\in\operatorname{Norm}_G(T)(S)$，\Term{lifting}{提升} \Term{weyl-group}{Weyl 群}的对称 $w_0$，\footnote{编者注：对称 $w_0$ 定义于 XXI 3.6.14。}即满足 $\operatorname{int}(n_0)U=U^-$（参见第 XXI 讲 3.6.14）；于是 $n_0^2\in T(S)$。说明 6.1.5 中的\Term{open-subset}{开集} $V_0=V_{n_0}$ 满足要求。事实上，
\[\begin{split}
U(S')U^-(S')
&=\operatorname{int}(n_0)\bigl(\operatorname{int}(n_0)^{-1}U(S')\cdot\operatorname{int}(n_0)^{-1}U^-(S')\bigr)\\
&=\operatorname{int}(n_0)\bigl(U^-(S')\cdot U(S')\bigr)\subset\operatorname{int}(n_0)\Omega'(S').
\end{split}\]
所以 $U(S')U^-(S')\cap V_0(S')\subset\operatorname{int}(n_0)\Omega'(S')\cap V_0(S')\subset\Omega'(S')$。
\end{proof}
\begin{SGAresult}{引理}{6.1.7}
考虑由第二\Term{projection}{投影}与 $T$ 到 $T/T'$ 的\Term{canonical}{典范}\Term{morphism}{态射}复合而得的\Term{morphism}{态射}
\[f:\Omega=U^-\cdot T\cdot U\longrightarrow T/T'.\]
则 $f$ “泛满足乘法性”（génériquement multiplicatif）：存在 $\Omega\times_S\Omega$ 的一个包含单位\Term{section}{截面}的\Term{open-subset}{开集} $V$（故由第 XVIII 讲 1.3，它相对\Term{schematically-dense}{概形稠密}），使对每个 $S'\to S$ 及 $(x,y)\in V(S')$，均有 $xy\in\Omega(S')$ 且 $f(xy)=f(x)f(y)$。
\end{SGAresult}
\begin{proof}
设 $x,y$ 为 $\Omega$ 在 $S'$ 上的两个\Term{section}{截面}，写成
\[x=utv,\qquad y=u't'v',\qquad
u,u'\in U^-(S'),\quad t,t'\in T(S'),\quad v,v'\in U(S').\]
令 $V_0$ 为 6.1.6 的\Term{open-subset}{开集}，$V$ 为条件“$vu'\in V_0(S')$”定义的 $\Omega\times_S\Omega$ 的\Term{open-subset}{开集}（即 $V_0$ 在逐点写成 $(x,y)\mapsto vu'$ 的\Term{morphism}{态射}下的\Term{inverse-image}{逆像}）。则 $V$ 满足要求。事实上，对 $(x,y)\in V(S')$，有
\[xy=(utv)(u't'v')=(ut)(vu')(t'v').\]
但 $vu'\in\Omega'(S')$，故
\[xy\in U^-(S')t\Omega'(S')t'U(S')\subset U^-(S')tt'T'(S')U(S');\]
这说明 $xy\in\Omega(S')$，且 $f(xy)=f(tt')=f(t)f(t')=f(x)f(y)$。
\end{proof}
\begin{SGAresult}{命题}{6.1.8}
存在一个群\Term{morphism}{态射} $f:G\to T/T'$，在 $T$ 上诱导\Term{canonical-projection}{典范投影}。$f$ 的核 $\operatorname{Ker}(f)$ 是 $G$ 的一个闭群子\Term{scheme}{概形}，在 $S$ 上\Term{smooth}{光滑}，\Term{fiber}{纤维}\Term{connected}{连通}。从 $G$ 到一个对 (fppf) \Term{separated}{分离}的 $S$ 上\Term{abelian-group}{交换群}值\Term{presheaf}{预层}的任意群\Term{morphism}{态射}都在 $\operatorname{Ker}(f)$ 上为零。
\end{SGAresult}
\SourcePage{171}
\begin{proof}
第一个断言立即由 4.1.11 得到。立即有 $\operatorname{Ker}(f)\cap\Omega=\Omega'$，这证明 $\operatorname{Ker}(f)$ 在单位\Term{section}{截面}的每一点处在 $S$ 上\Term{smooth}{光滑}。\footnote{编者注：已加入“在\Term{unit-section}{单位截面}的每一点处”，并在证明末尾加入对 VI$_B$ 3.10 的引用。此外，下一句的“\Term{prescheme}{预概形}”改为“\Term{presheaf}{预层}”。}由 5.6.9 (ii)，从 $G$ 到一个对 (fppf) \Term{separated}{分离}的\Term{abelian-group}{交换群}值\Term{presheaf}{预层}的任意\Term{morphism}{态射} $\varphi$，都在 $U,U^-$ 上为零。由第 XX 讲 2.7，$\varphi$ 也在 $T'$ 上为零，故在 $\Omega'$ 上为零。沿用 5.7.10 的记号，\Term{monoid}{幺半群} $U_1$ 包含于 $\operatorname{Ker}(f)(S)$，故
\[\operatorname{Ker}(f)=\bigcup_{u\in U_1}u\Omega'.\]
一方面，任意上述 $\varphi$ 均在 $\operatorname{Ker}(f)$ 上为零；另一方面，$\operatorname{Ker}(f)$ 的\Term{fiber}{纤维}\Term{connected}{连通}，故由第 VI$_B$ 讲 3.10，它在 $S$ 上\Term{smooth}{光滑}。
\end{proof}
\SGAsection{6.2}{约化群的导群}
\begin{SGAresult}{定理}{6.2.1}
设 $S$ 为\Term{scheme}{概形}，$G$ 为 $S$ 上的\Term{reductive-group}{约化群}。

(i) $D_S(G)=\operatorname{Hom}_{S\text{-gr.}}(G,\mathbb G_{m,S})$ 可由一个 $S$ 上的\Term{twisted-constant-group}{扭常值群}表示，其在 $s\in S$ 处的类型为 $\mathbb Z^{\operatorname{rgred}(G_s)-\operatorname{rgss}(G_s)}$。

(ii) 记 $\operatorname{corad}(G)=D_S(D_S(G))$，它因而是一个 $S$ 上的\Term{torus}{环面}。\Term{biduality}{双对偶}\Term{morphism}{态射}（参见第 VIII 讲第 1 节）$f_0:G\to\operatorname{corad}(G)$ \Term{smooth}{光滑}且\Term{surjectivity}{满射}。

(iii) 复合\Term{morphism}{态射} $\operatorname{rad}(G)\to G\to\operatorname{corad}(G)$ 是\Term{isogeny}{同源}（参见 4.2.9）。

(iv) $f_0$ 的核，记作 $\operatorname{d\text{é}r}(G)=\operatorname{Ker}(f_0)$，是 $G$ 的一个闭群子\Term{scheme}{概形}，在 $S$ 上\Term{semisimple}{半单}，称为 $G$ 的\Term{derived-group}{导群}。若 $G$ \Term{semisimple}{半单}，则 $\operatorname{d\text{é}r}(G)=G$。

(v) 从 $G$ 到一个对 (fppf) \Term{separated}{分离}的 $S$ 上\Term{abelian-group}{交换群}值\Term{presheaf}{预层}的任意群\Term{morphism}{态射}都在 $\operatorname{d\text{é}r}(G)$ 上为零，故经由 $f_0$ \Term{resolution}{分解}。
\end{SGAresult}
\begin{proof}
本定理的所有断言对\Term{etale}{平展}\Term{topology}{拓扑}都是\Term{local-scheme}{局部}的，故可归结到 $G$ 在 $S$ 上\Term{split-group}{分裂}的情形。考虑 6.1.8 的\Term{morphism}{态射} $f$。由 6.1.8 的最后一个断言，立即有\Term{isomorphism}{同构}
\[\operatorname{Hom}_{S\text{-gr.}}(G,\mathbb G_{m,S})\xrightarrow\sim
\operatorname{Hom}_{S\text{-gr.}}(T/T',\mathbb G_{m,S}),\]
\SourcePage{172}
这证明了 (i)，继而证明 (ii)，并给出\Term{commutative-diagram}{交换图}
\[\begin{tikzcd}
G\arrow[r,"f_0"]\arrow[dr,"f"']&\operatorname{corad}(G)\arrow[d,"\sim"]\\
&T/T'.
\end{tikzcd}\]
由 6.1.8 得到 (v)，由 6.1.2 得到 (iii)。又 $\operatorname{Ker}(f)=\operatorname{Ker}(f_0)$，由 6.1.8 推出 $\operatorname{d\text{é}r}(G)$ 在 $S$ 上\Term{smooth}{光滑}，\Term{fiber}{纤维}\Term{connected}{连通}；只须验证其\Term{fiber}{纤维}\Term{semisimple}{半单}。它们作为\Term{reductive-group}{约化群}的\Term{normal-subgroup}{正规子群}，由第 XIX 讲 1.7 是\Term{reduced}{约化}的。由 (iii)，$\operatorname{rad}(G)\cap\operatorname{d\text{é}r}(G)$ \Term{finite}{有限}，从而 $\operatorname{d\text{é}r}(G)$ 的\Term{fiber}{纤维}确实\Term{semisimple}{半单}。
\end{proof}
\textbf{注记 6.2.2。} a) 由构造，若 $G$ \Term{split-group}{分裂}，则 $\operatorname{d\text{é}r}(G)$ 是 $G$ 中由 $U_\alpha$（$\alpha\in R$）生成的 (fppf) 子\Term{sheaf}{层}。（甚至只须取 $\alpha\in\Delta\cup-\Delta$ 的 $U_\alpha$，其中 $\Delta$ 为 $R$ 的一组基。）

b)\footnote{编者注：下文详细补充了原版，并删去“$\operatorname{d\text{é}r}(G)$ 是 $G$ 的对 (fppf) \Term{separated}{分离}的\Term{commutator-presheaf}{换位子预层}”这一断言。} 令 $C$ 为 $G$ 的\Term{commutator-presheaf}{换位子预层}，即向每个 $S'\to S$ 赋予 $G(S')$ 的换位\Term{subgroup}{子群}的 $S$ 上群值\Term{functor}{函子}（该群即 $G(S')$ 中由 $xyx^{-1}y^{-1}$、$x,y\in G(S')$ 生成的\Term{subgroup}{子群}），令 $\widetilde C$ 为其相伴的 (fppf) \Term{sheaf}{层}。因商 $G/\operatorname{d\text{é}r}(G)=T/T'$ 交换，$\operatorname{d\text{é}r}(G)$ 包含 $C$，故也包含 $\widetilde C$（参见第 IV 讲 4.3.12）。
另一方面，商\Term{presheaf}{预层} $G/\widetilde C$ \Term{separated}{分离}（第 IV 讲 4.4.8.1），由 (v)，$\operatorname{d\text{é}r}(G)\subset\widetilde C$，所以 $\operatorname{d\text{é}r}(G)=\widetilde C$，即 $\operatorname{d\text{é}r}(G)$ 是 $G$ 的 (fppf) \Term{commutator-sheaf}{换位子层}。
最后注意，$C$ 作为 $G$ 的一个子\Term{presheaf}{预层}，是\Term{separated}{分离}的，但一般不等于 $\operatorname{d\text{é}r}(G)$：例如 $\operatorname{d\text{é}r}(\operatorname{SL}_2)=\operatorname{SL}_2$，但 $\operatorname{SL}_2(\mathbb F_2)\simeq\mathfrak S_3$ 不等于其\Term{derived-group}{导群}。

c) 若 $S$ 是\Term{algebraically-closed-field}{代数闭域} $k$ 的谱，则 $\operatorname{d\text{é}r}(G)(k)$ 是 $G(k)$ 的换位\Term{subgroup}{子群}（第 VI$_B$ 讲 7.10）。

\textbf{6.2.3。} 现在考虑两个\Term{exact-sequence}{正合列}
\[1\longrightarrow\operatorname{rad}(G)\longrightarrow G\longrightarrow\operatorname{ss}(G)\longrightarrow1,\]
\[1\longrightarrow\operatorname{d\text{é}r}(G)\longrightarrow G\longrightarrow\operatorname{corad}(G)\longrightarrow1.\]
因 $\operatorname{rad}(G)$ 在 $G$ 中为\Term{central-subgroup}{中心子群}，$G$ 中乘法定义一个群\Term{morphism}{态射}
\[u:\operatorname{rad}(G)\times_S\operatorname{d\text{é}r}(G)\longrightarrow G;\]
由 6.2.1 (iii)，它为\Term{covering}{覆盖}，故\Term{surjectivity}{满射}且\Term{flat}{平坦}（第 VI$_B$ 讲 9.2 (xi)）。\footnote{编者注：事实上，$G$ 是 $\operatorname{rad}(G)\times_S\operatorname{d\text{é}r}(G)$ 模 $\operatorname{Ker}(u)$ 的 (fppf) 商，而 $\operatorname{Ker}(u)$ 为\Term{multiplicative-type}{乘法型}群，故在 $S$ 上\Term{flat}{平坦}。所以由 VI$_B$ 9.2 (xi)，\Term{morphism}{态射} $u$ \Term{flat}{平坦}。}其核与 $\operatorname{rad}(G)\cap\operatorname{d\text{é}r}(G)$ \Term{isomorphism}{同构}，后者也是 $\operatorname{rad}(G)\to\operatorname{corad}(G)$ 的核，故为 $\operatorname{rad}(G)$ 的一个\Term{finite}{有限}\Term{multiplicative-type}{乘法型}\Term{subgroup}{子群}。
\SourcePage{173}
对\Term{morphism}{态射} $G\to\operatorname{corad}(G)\times_S\operatorname{ss}(G)$ 作同样论证，其核为 $\operatorname{d\text{é}r}(G)\cap\operatorname{rad}(G)$。因而有：
\begin{SGAresult}{命题}{6.2.4}
设 $G$ 为 $S$ 上的\Term{reductive-group}{约化群}。\Term{morphism}{态射}
\[\operatorname{rad}(G)\times_S\operatorname{d\text{é}r}(G)\longrightarrow G,\qquad
G\longrightarrow\operatorname{corad}(G)\times_S\operatorname{ss}(G),\qquad
\operatorname{rad}(G)\longrightarrow\operatorname{corad}(G)\]
均为\Term{central-isogeny}{中心同源}，且它们的核相互\Term{isomorphism}{同构}。
\end{SGAresult}
\begin{SGAresult}{推论}{6.2.5}
下列条件\Term{equivalence}{等价}：

(i) $G$ 是一个\Term{semisimple-group}{半单群}与一个\Term{torus}{环面}的乘积。

(ii) $\operatorname{rad}(G)\times_S\operatorname{d\text{é}r}(G)\xrightarrow\sim G$。

(iii) $G\xrightarrow\sim\operatorname{corad}(G)\times_S\operatorname{ss}(G)$。

(iv) $\operatorname{rad}(G)\cap\operatorname{d\text{é}r}(G)=e$。
\end{SGAresult}
\textbf{6.2.6。} 暂时回到\Term{split-group}{分裂群}的情形，沿用 6.1 的记号。置 $N=M\cap\mathscr V(R^*)^\perp$，故 $T'=D_S(M/N)$。已见 $U^-\cdot T'\cdot U$ 是 $\operatorname{d\text{é}r}(G)$ 中单位\Term{section}{截面}的一个\Term{open-neighborhood}{开邻域}。因此
\[\operatorname{Lie}(\operatorname{d\text{é}r}(G)/S)=\mathfrak t'\oplus\bigoplus_{\alpha\in R}\mathfrak g^\alpha.\]
因 $\alpha\in R$ 在 $T'$ 上诱导的\Term{character}{特征标}均非零且互异（参见第 XXI 讲 1.2.5；其实在 6.1.2 中已用到此事实），$R$ 是 $G$ 关于 $T$ 的一个\Term{root-system}{根系}。于是（因 $U_\alpha\subset\operatorname{d\text{é}r}(G)$）立即可见，$\operatorname{d\text{é}r}(G)$ 的指数\Term{morphism}{态射}“就是”$G$ 的指数\Term{morphism}{态射}，\Term{coroot}{余根}也一样。
由此得到：
\begin{SGAresult}{命题}{6.2.7}
沿用上述记号，$(\operatorname{d\text{é}r}(G),T',M/N,R)$ 是\Term{root-datum}{根数据}为 $\operatorname{d\text{é}r}(\mathcal R(G))$ 的\Term{split-group}{分裂群}。\Term{canonical}{典范}\Term{morphism}{态射} $\operatorname{d\text{é}r}(G)\to G$ 通过\Term{functor}{函子}性给出第 XXI 讲 6.5 的\Term{canonical}{典范}\Term{root-datum}{根数据}\Term{morphism}{态射} $\mathcal R(G)\to\operatorname{d\text{é}r}(\mathcal R(G))$。
\end{SGAresult}
\textbf{附注。} 读者可作为练习，构造与第 XXI 讲 6.5.7 的\Term{root-datum}{根数据}图最左三列对应的\Term{split-group}{分裂群}图。
\begin{SGAresult}{命题}{6.2.8}
设 $S$ 为\Term{scheme}{概形}，$G$ 为 $S$ 上的\Term{reductive-group}{约化群}，$\operatorname{d\text{é}r}(G)$ 为其\Term{derived-group}{导群}。

(i) 对 $G$ 的每个\Term{maximal-torus}{极大环面} $T$，$T\cap\operatorname{d\text{é}r}(G)$ 是 $\operatorname{d\text{é}r}(G)$ 的一个\Term{maximal-torus}{极大环面}。对 $\operatorname{d\text{é}r}(G)$ 的每个\Term{maximal-torus}{极大环面} $T'$，$\operatorname{Centr}_G(T')=\operatorname{rad}(G)\cdot T'$ 是 $G$ 的一个\Term{maximal-torus}{极大环面}。上述两个构造互逆，在 $G$ 与 $\operatorname{d\text{é}r}(G)$ 的\Term{maximal-torus}{极大环面}之间建立\Term{bijection}{双射}。
\SourcePage{174}
(ii) 对 $G$ 的每个\Term{borel-subgroup}{Borel 子群} $B$，$B\cap\operatorname{d\text{é}r}(G)$ 是 $\operatorname{d\text{é}r}(G)$ 的一个\Term{borel-subgroup}{Borel 子群} $B'$。有 $B^{\prime u}=B^u$。对 $\operatorname{d\text{é}r}(G)$ 的每个\Term{borel-subgroup}{Borel 子群} $B'$，$\operatorname{Norm}_G(B')=\operatorname{rad}(G)\cdot B'$ 是 $G$ 的一个\Term{borel-subgroup}{Borel 子群}。上述映射互逆，在 $G$ 与 $\operatorname{d\text{é}r}(G)$ 的\Term{borel-subgroup}{Borel 子群}之间建立\Term{bijection}{双射}。
\end{SGAresult}
\begin{proof}
由\Term{maximal-torus}{极大环面}的\Term{local-scheme}{局部}\Term{conjugacy}{共轭}定理及\Term{derived-group}{导群}的构造，(i) 中只须再证明：若 $T$ 为 $G$ 的一个\Term{maximal-torus}{极大环面}，则
\[T=(T\cap\operatorname{d\text{é}r}(G))\cdot\operatorname{rad}(G)
=\operatorname{Centr}_G(T\cap\operatorname{d\text{é}r}(G)).\]
第一个等式显然（因为可归结为\Term{splitting}{分裂}情形）；第二个立即由此得到，因为 $\operatorname{rad}(G)$ 在 $G$ 中为\Term{central-subgroup}{中心子群}，故 $T=\operatorname{Centr}_G(T)=\operatorname{Centr}_G(T\cap\operatorname{d\text{é}r}(G))$。对 (ii) 作同样论证。
\end{proof}
\SGAsection{6.3}{具有交换商的子群}
\textbf{6.3.1。} 设 $G$ 为 $S$ 上的\Term{reductive-group}{约化群}。若 $H$ 为 $G$ 的一个群子\Term{sheaf}{层}，则下列条件\Term{equivalence}{等价}：

-- $H$ 包含 $\operatorname{d\text{é}r}(G)$。

-- $H$ \Term{normal}{正规}，且 $G/H$ 交换。

此时，\Term{canonical}{典范}\Term{morphism}{态射} $f_0:G\to\operatorname{corad}(G)$ 将 $H$ 映到 $\operatorname{corad}(G)$ 的一个子\Term{sheaf}{层} $f_0(H)$；有
\[\begin{aligned}
G/H&\simeq\operatorname{corad}(G)/f_0(H),&
H/\operatorname{d\text{é}r}(H)&\simeq f_0(H),\\
\operatorname{d\text{é}r}(G)&=\operatorname{d\text{é}r}(H),&
H&=f_0^{-1}(f_0(H)).
\end{aligned}\]
因 $\operatorname{d\text{é}r}(G)$ 在 $S$ 上\Term{smooth}{光滑}，\Term{fiber}{纤维}\Term{connected}{连通}，\footnote{编者注：下文详细补充了原版。特别是加入了原版中隐含的结论：$G$ 的每个\Term{smooth}{光滑}、\Term{fiber}{纤维}\Term{connected}{连通}且包含 $\operatorname{d\text{é}r}(G)$ 的\Term{subgroup}{子群}在 $G$ 中均闭。}由第 IV 讲 5.3.1、6.3.1 及第 IV$_B$ 讲 9.2，映射 $H\mapsto f_0(H)$ 在 $G$ 中包含 $\operatorname{d\text{é}r}(G)$、在 $S$ 上\Term{smooth}{光滑}且\Term{fiber}{纤维}\Term{connected}{连通}的群子\Term{scheme}{概形}（分别为闭群子\Term{scheme}{概形}）与 $\operatorname{corad}(G)$ 中在 $S$ 上\Term{smooth}{光滑}且\Term{fiber}{纤维}\Term{connected}{连通}的群子\Term{scheme}{概形}（分别为闭群子\Term{scheme}{概形}）之间建立\Term{bijection}{双射}。
若 $H'$ 是 $\operatorname{corad}(G)$ 的一个在 $S$ 上\Term{smooth}{光滑}且\Term{fiber}{纤维}\Term{connected}{连通}的群子\Term{scheme}{概形}，则 $H'$ 在 $S$ 上\Term{finite-presentation}{有限呈示}（第 VI$_B$ 讲 5.5），且其\Term{fiber}{纤维}为\Term{torus}{环面}（因为 $\operatorname{corad}(G)$ 的\Term{fiber}{纤维}是\Term{torus}{环面}），故由第 X 讲 8.2，$H'$ 是 $\operatorname{corad}(G)$ 的一个\Term{subtorus}{子环面}，因而在 $\operatorname{corad}(G)$ 中闭（第 IX 讲 2.6）。因此，$G$ 的每个\Term{smooth}{光滑}、\Term{fiber}{纤维}\Term{connected}{连通}且包含 $\operatorname{d\text{é}r}(G)$ 的\Term{subgroup}{子群}，在 $G$ 中均闭（见编者脚注 60）。

\textbf{6.3.2。} 若 $H$ 是 $G$ 的一个在 $S$ 上\Term{smooth}{光滑}、\Term{fiber}{纤维}\Term{connected}{连通}、在 $G$ 中\Term{normal}{正规的}闭群子\Term{scheme}{概形}，则 $H$ \Term{reductive-group}{约化}。若还有 $H\supset\operatorname{d\text{é}r}(G)$，则 $\operatorname{d\text{é}r}(H)=\operatorname{d\text{é}r}(G)$，且 $f_0(H)$ 与 $\operatorname{corad}(H)$ 等同。于是证明了：
\SourcePage{175}
\begin{SGAresult}{命题}{6.3.3}
设 $G$ 为 $S$ 上的\Term{reductive-group}{约化群}。$G$ 的每个在 $G$ 中\Term{normal}{正规}、商交换（即包含 $\operatorname{d\text{é}r}(G)$）、在 $S$ 上\Term{smooth}{光滑}且\Term{fiber}{纤维}\Term{connected}{连通}的群子\Term{scheme}{概形} $H$，\footnote{编者注：已删去 $H$ 在 $G$ 中\Term{retrocompact}{逆紧}（rétrocompact）的假设，该假设自动成立：由 VI$_B$ 5.5，$G,H$ 在 $S$ 上\Term{separated}{分离}且\Term{quasi-compact}{拟紧}，故由 EGA IV$_1$，1.1.2 (v)，$H\hookrightarrow G$ \Term{quasi-compact}{拟紧}。}均闭且\Term{reductive-group}{约化}。有 $\operatorname{d\text{é}r}(H)=\operatorname{d\text{é}r}(G)$，且 $f_0(H)$ 与 $\operatorname{corad}(H)$ 等同；有
\[G/H\simeq\operatorname{corad}(G)/\operatorname{corad}(H),\qquad
H=(H\cap\operatorname{rad}(G))\cdot\operatorname{d\text{é}r}(G).\]
此外，$H\mapsto f_0(H)$ 在具有上述性质的 $G$ 的\Term{subgroup}{子群} $H$ 的\Term{set}{集合}与 $\operatorname{corad}(H)$ 的\Term{subtorus}{子环面}的\Term{set}{集合}之间定义\Term{bijection}{双射}。
\end{SGAresult}
再次应用 Noether 定理（第 IV 讲 5.3.1 和 6.3.1），得到：
\begin{SGAresult}{命题}{6.3.4}
设 $S$ 为\Term{scheme}{概形}，$G$ 为 $S$ 上的\Term{reductive-group}{约化群}，$T$ 为 $G$ 的一个\Term{maximal-torus}{极大环面}。对 $G$ 的每个上述\Term{subgroup}{子群} $H$，$T\cap H$ 是 $G$ 的一个\Term{maximal-torus}{极大环面}，且
\[G/H\simeq T/(T\cap H),\qquad H=(T\cap H)\cdot\operatorname{d\text{é}r}(G).\]
此外，$H\mapsto T\cap H$ 在 $G$ 的上述\Term{subgroup}{子群} $H$ 的\Term{set}{集合}与 $T$ 中包含 $T\cap\operatorname{d\text{é}r}(G)$ 的\Term{subtorus}{子环面}的\Term{set}{集合}之间给出\Term{bijection}{双射}。
\end{SGAresult}
\SGAsection{}{参考文献}
\begin{description}
\item[{[Ch05]}] C. Chevalley，\emph{\Term{semisimple}{半单}\Term{algebraic-group}{代数群}的分类（Classification des groupes algébriques semi-simples）}（P. Cartier、A. Grothendieck、M. Lazard 协作），Collected Works，第 3 卷，Springer，2005 年。
\item[{[DG70]}] M. Demazure、P. Gabriel，\emph{\Term{algebraic-group}{代数群}（Groupes algébriques）}，Masson \& North-Holland，1970 年。
\item[{[Gi71]}] J. Giraud，\emph{\Term{nonabelian-cohomology}{非阿贝尔上同调}（Cohomologie non abélienne）}，Springer-Verlag，1971 年。
\item[{[Se64]}] J.-P. Serre，\emph{Galois \Term{cohomology}{上同调}（Cohomologie galoisienne）}，Springer-Verlag，1964 年；第 5 版，1994 年。
\end{description}
\SourcePage{176}
% 原 PDF 第 176 页为空白。
